\documentclass[11pt]{amsart}
\usepackage{lmodern}
\usepackage[expansion=false]{microtype}
\usepackage[a4paper,textwidth=152mm,textheight=242mm,centering]{geometry}
\usepackage{amsmath,amssymb,mathtools,booktabs,enumitem}
\usepackage[colorlinks,linkcolor=blue!50!black,citecolor=blue!50!black,urlcolor=blue!50!black]{hyperref}
\usepackage{xcolor}
\setlength{\emergencystretch}{1.4em}
\numberwithin{equation}{section}
\newtheorem{theorem}{Theorem}[section]
\newtheorem{lemma}[theorem]{Lemma}
\newtheorem{proposition}[theorem]{Proposition}
\newtheorem{corollary}[theorem]{Corollary}
\theoremstyle{definition}
\newtheorem{definition}[theorem]{Definition}
\newtheorem{example}[theorem]{Example}
\newtheorem{remark}[theorem]{Remark}
\theoremstyle{plain}
\newtheorem{mainthm}{Theorem}
\newtheorem{maincor}[mainthm]{Corollary}
\renewcommand{\themainthm}{\Alph{mainthm}}
\DeclareMathOperator{\lcm}{lcm}
\DeclareMathOperator{\Tr}{Tr}
\DeclareMathOperator{\FPdim}{FPdim}
\DeclareMathOperator{\Irr}{Irr}
\DeclareMathOperator{\Gal}{Gal}
\newcommand{\QQ}{\mathbb Q}
\newcommand{\ZZ}{\mathbb Z}
\newcommand{\C}{\mathcal C}
\allowdisplaybreaks[2]

\title[A rank bound for simple integral modular data]{A rank bound for simple integral modular data}
\author{Sebastian Burciu}
\address[S. Burciu]{Simion Stoilow Institute of Mathematics of the Romanian Academy, Bucharest, Romania}
\email{sebastian.burciu@imar.ro}
\author{Jingcheng Dong}
\address[J. Dong]{College of Mathematics and Statistics, Nanjing University of Information Science and Technology, Nanjing, China}
\email{jcdong@nuist.edu.cn}
\author{S\'ebastien Palcoux}
\address[S. Palcoux]{Beijing Institute of Mathematical Sciences and Applications, Beijing, China}
\email{sebastien.palcoux@gmail.com}
\date{}
\hypersetup{pdftitle={A rank bound for simple integral modular data},pdfauthor={Sebastian Burciu, Jingcheng Dong and S\'ebastien Palcoux}}
\begin{document}
\begin{abstract}
We prove that a non-pointed based-simple commutative integral fusion ring with a symmetric algebraic Drinfeld map and rotation-admissible cyclotomic modular data has rank at least $45$. Character orthogonality and the primary components of the twists bound the multiplicities of the largest degrees. An exact enumeration excludes the remaining degree lists.
\end{abstract}
\maketitle

\section{The rank problem}
It remains open whether there exists a non-pointed simple integral modular fusion category \cite[Question~1.3 and Theorem~5.8]{LPR}. Burciu, Dong and Palcoux proved that such a category has rank at least $32$ \cite[Theorem~22.2]{BDPcore}. We give an algebraic bound in terms of the fusion ring and its modular data.

\smallskip\noindent
A commutative fusion ring $R$ is a commutative unital ring with finite
basis $(x_i)_{i\in I}$ over $\ZZ$, unit $x_0$ (the vacuum), nonnegative integral
structure constants $N_{ij}^k$, and a basis-preserving ring involution
$x_i\mapsto x_{i^*}$ such that $N_{ij}^0=\delta_{i,j^*}$.
A character is a unital homomorphism $R\to\mathbb C$. Write $d_i$ for the value at $x_i$ of the unique positive character,
$r=|I|$ for its rank, and $D=\sum_i d_i^2$ for its global dimension.
The ring is \emph{integral} if all $d_i$ are integers, \emph{pointed}
if all equal one, and \emph{based-simple} if every nonunit $x_i$
generates the full basis under duality and passage to constituents of
products. A \emph{symmetric algebraic Drinfeld map} is a bijection
from the probability basis $c_i=x_i/d_i$ to the complex characters
$\phi_i$ satisfying $\phi_i(c_j)=\phi_j(c_i)$. Put
$\chi_i(j)=\phi_i(x_j)$, $S_{ij}=d_i\chi_i(j)$, and $s=S/\sqrt D$.

\begin{mainthm}\label{alg:rank}
Let $R$ be a non-pointed based-simple commutative integral fusion ring whose probability basis admits a symmetric algebraic Drinfeld map. Suppose that $R$ carries rotation-admissible cyclotomic modular data as in Definition~\ref{alg:mod-data}. Then $\operatorname{rk}(R)\ge45$.
\end{mainthm}

\begin{definition}\label{alg:mod-data}
\emph{Cyclotomic modular data} consist of roots of unity $\theta_i$
with $\theta_0=1$ and a root of unity $\gamma$ of two-power order.
Set $T=\operatorname{diag}(\theta_i)$,
$N=\operatorname{ord}(T)$ and $L=\operatorname{ord}(T/\gamma)$, and write $\zeta_n=e^{2\pi i/n}$.
The primes dividing $N$ must be precisely those dividing $D$, and
the assignments
\[
 W=\left(\begin{smallmatrix}0&-1\\1&0\end{smallmatrix}\right)
 \longmapsto s,\qquad
 U=\left(\begin{smallmatrix}1&1\\0&1\end{smallmatrix}\right)
 \longmapsto T/\gamma
\]
must define a representation $\rho$ of $\mathrm{SL}_2(\ZZ/L)$ over
$\QQ(\zeta_L)$ with
\begin{equation}\label{alg:covariance}
 \sigma_a\rho(g)=\rho(D_a gD_a^{-1}),\qquad
 D_a=\operatorname{diag}(a,1),\quad \sigma_a(\zeta_L)=\zeta_L^a
 \quad (a\in(\ZZ/L)^\times).
\end{equation}
For $m\ge1$ define the linear functional
\begin{equation}\label{alg:rotation}
 L_m^i(x_c)=\sum_{a,b}s_{ia}s_{ib}(\theta_a/\theta_b)^m N_{ab}^c .
\end{equation}
A \emph{rotation certificate} for $(i,m,c)$ is a finite multiset
of $m$th roots of unity whose $t$th powers sum to $L_{m/t}^i(x_c^t)$
for each positive divisor $t$ of $m$. We call the data
\emph{rotation-admissible} if every $(i,m,c)$ has a certificate.
\end{definition}

Equip an integral modular fusion category with its canonical positive spherical structure \cite[Propositions~8.23--8.24]{ENOannals}. Its normalized $S$-matrix gives the character pairing, and generalized Frobenius--Schur rotations give the required power traces \cite{DLNCongruence,NSind}.

Write $D=qM^2$ with $q$ square-free. Integrality of the formal codegrees gives integers $b_i=M/d_i$ satisfying $\sum_i b_i^{-2}=q$ \cite[Lemma~2.3]{HR}, \cite[\S4]{ABPP}. Small rank forces some $b_i$ to be small. Reduction of the characters at each prime dividing $D$ partitions the labels into blocks. Orthogonality within each block and the primary components of the twists restrict the multiplicities and character rows of the largest degrees. These restrictions give a finite enumeration in Section~\ref{sec:verify}.

\section{Character arithmetic}
Throughout the proof, $R$ and its modular data satisfy the hypotheses of Theorem~\ref{alg:rank}.

\begin{lemma}[The normalized character table]
The matrix $s$ is symmetric and unitary, $s^2$ is the duality
permutation $C$, and
\begin{equation}
 \sum_j\chi_i(j)\overline{\chi_k(j)}=\delta_{ik}f_i,\qquad
 f_i=D/d_i^2\in\ZZ,\qquad S_{ij}/\lcm(d_i,d_j)\in\overline{\mathbb Z} .
\end{equation}
The character field $K_i=\QQ(\chi_i(j):j\in I)$ has conductor dividing
$f_i$. Galois permutes the labels, preserving degree, with
$\sigma S_{ij}=S_{\sigma i,j}=S_{i,\sigma j}$.
\end{lemma}
\begin{proof}
The fusion matrices are commuting normal matrices, their adjoints
being given by duality. Their distinct character rows are orthogonal.
Symmetry gives $SS^*=S^*S=\operatorname{diag}(d_i^2f_i)$, where
$f_i=\sum_j|\chi_i(j)|^2$. This diagonal commutes with $S$ and its
vacuum column has no zero, so it is $DI$. Also $\bar S=CS$.
The entries $\chi_i(j)$ are algebraic integers. This gives integrality
of the rational $f_i$, and B\'ezout applied to $S_{ij}/d_i,S_{ij}/d_j$
gives the lcm assertion. Galois permutes the characters. Their
codegrees, hence their degrees, are fixed. Symmetry then gives
covariance in both indices. The row-conductor theorem
\cite[Theorem~2.9]{BDPcore} applies to these symmetric integral data.
\end{proof}

As $\rho(-I)=C$, dual labels have the same twist. The identity
$WUW=U^{-1}WU^{-1}$ and its vacuum entry give
\[
 sTs=\xi T^{-1}sT^{-1},\quad \xi=\gamma^3=p^+/\sqrt D,\quad p^+=\sum_i d_i^2\theta_i .
\]
For $v=(d_i\theta_i)_i$, its vacuum column gives
$(s^{-1}v)_l=\sqrt D\,\xi s_{0l}/\theta_l$. Diagonalizing fusion
and using the inverse of the same identity gives the balancing identity
\[
 \sum_k N_{i^*j}^k d_k\theta_k
 =\sqrt D\,\xi(sT^{-1}s^{-1})_{ji}=\theta_i\theta_j S_{ij}.
\]
Thus $S_{ij}\in\QQ(\zeta_N)$. Write $h_a=\rho(\operatorname{diag}(a,a^{-1}))$.
By \eqref{alg:covariance}, $h_as=\sigma_a(s)$. Row covariance makes
$h_a$ a signed permutation fixing the vacuum line. Conjugacy of $U$ to $U^{a^2}$ by
the diagonal matrix therefore gives $\gamma^{a^2-1}=1$ for every
unit $a$ modulo $L$. At a prime of $L$ outside $D$, both generating
unipotents of the local factor act by scalars.

\begin{lemma}[Arithmetic of rotation traces]\label{alg:rotation-arithmetic}
Every nonnegative integral combination $V$ of basis elements has
rotation certificates. All their power traces belong to $\QQ(\zeta_N)$.
For $p\mid m$ and $y$ with algebraic-integer coefficients,
$L_m^i(y)^p\equiv L_{m/p}^i(y^p)$ at every place over $p$.
For $m=N_{p'}:=N/p^{v_p(N)}$, $L_m^0(x_c)\in\ZZ$.
\end{lemma}
\begin{proof}
Take the summands of $V$ as letters, with distinct copies distinguished.
For each cyclic orbit of length-$m$ words of least period $b$, let $Y$
be the product along a period. Induce its certificate from $C_{m/b}$
to $C_m$. This exists by hypothesis for $b=1$ and by induction on $m$
otherwise. Its trace at $t\mid m$ is zero unless $b\mid t$, when it is
$bL_{m/t}^i(Y^{t/b})$. Summing counts each length-$t$ word once.
Divisor-power traces are in $\QQ(\zeta_N)$ by \eqref{alg:rotation}
and $s_{ia}s_{ib}=S_{ia}S_{ib}/D$. Other power traces are their
Galois conjugates. Reduction of the roots gives the congruence on the
basis, hence by Frobenius on all $y$. For the last assertion the
integral trace lies in both $\QQ(\zeta_m)$ and $\QQ(\zeta_{p^{v_p(N)}})$,
the latter by the rational vacuum weights in \eqref{alg:rotation}.
\end{proof}

\begin{lemma}[Vanishing and divisibility]\label{lem:basic-table}
Only the unit has degree one. The dimension is square-full, with
at least three prime divisors, and every nonunit degree has at least
two prime divisors. For nonunits $i,j$ one has $|S_{ij}|<d_id_j$,
with $S_{ij}=0$ when $(d_i,d_j)=1$. After deleting any single nonunit
Galois orbit the remaining nonunit degrees have gcd one.
If $D$ is odd, every nonunit has a distinct dual.
\end{lemma}
\begin{proof}
The invertibles span a proper based subring, so reduce to the unit.
This property is called \emph{perfection}. By Burnside vanishing \cite[Theorem~1.2]{BPBurnside}, the product
$\prod_i x_i x_{i^*}$ is supported only on the degree character.
Its integral unit coefficient is $\prod_i d_i^2/D$. Since $d_i^2\mid D$,
every prime divisor of $D$ occurs at least twice. For fixed nonunit $i$ the equality
labels $|S_{ij}|=d_id_j$ span a based subring
\cite[Lemma~2.35]{BPBurnside}. Since $f_i<D$ it is proper.
For coprime degrees every conjugate of the algebraic integer
$S_{ij}/(d_id_j)$ therefore has modulus less than one. Vacuum
orthogonality gives $\sum_{j\ne0}d_j\chi_i(j)=-1$. If $d_i$ were a
prime power, every nonzero term would be divisible by that prime.
The same identity gives gcd one for all nonunit degrees, hence at
least three dimension primes. By symmetry every nonunit column has a zero. A row vanishing
there vanishes on its entire Galois orbit, giving the deletion claim.

For odd $D$, the unit coefficient
$\tau(y)=D^{-1}\sum_i d_i^2\phi_i(y)$ reduces above two to the sum
of the characters. Thus $\tau(x_i^2)\equiv\tau(x_i)^2\pmod2$,
or $\delta_{i,i^*}=\delta_{i,0}$.
\end{proof}

\paragraph{Ordered denominators.}
Write $D=qM^2$ with $q$ square-free. Then $d_i\mid M$. Put
$b_i=M/d_i$ and order the nonunit denominators increasingly, with
repetition. For a prime $p$ write $M_p=p^{v_p(M)}$. A $p$-free
degree is prime to $p$, and a pure $p$-denominator is a power of $p$.
Primes dividing $q$ are called ramified, and the other primes of $D$
unramified. Dividing the mass $\sum_{i\in J}d_i^2$ of a set of
labels by $M^2$ gives its reciprocal-square mass $\sum_{i\in J}b_i^{-2}$.
Thus
\begin{equation}
 q=M^{-2}+\sum_{i=1}^{r-1}b_i^{-2},\qquad qb_1^2<r .
\end{equation}
After a prefix of length $j\ge1$ write $L=\lcm(b_1,\ldots,b_j)$,
$E=L^2(q-\sum_{i\le j}b_i^{-2})$, $h=r-1-j$. For the next $b$,
\begin{equation}\label{eq:interval}
 b\ge b_j,\qquad L^2<Eb^2\le(h+1/36)L^2 .
\end{equation}
Indeed $M/b\ge6$. For $g=(L,b)$ the new pair is
$((b/g)L,(b/g)^2 E-(L/g)^2)$. With two degrees $d\le e$ remaining,
write $M=kL$. Each prefix degree is a multiple of $k$. Deleting
the orbit of either remaining label shows $(k,de)=1$. Thus
$d,e\mid L$, $Ek^2=1+d^2+e^2$, and $kL/e\ge b_j$.
Reduction modulo four and five gives $(k,10)=1$ and then
$E\not\equiv0,4,7\pmod8$. With one degree left, deletion gives
$k=1$, $E=1+d^2$.

For mutually orthogonal character rows $v$ perpendicular to the degree
row $d$, regarded as columns, orthogonal projection gives
\begin{equation}
 I-\sum_v vv^*/\|v\|^2\succeq dd^*/D .
\end{equation}
This inequality remains valid on every set of coordinates.

\paragraph{Cyclotomic estimates.}
Write $\langle z\rangle$ for the mean of the Galois conjugates of $z$
and $C(z)=\langle |z|^2\rangle$. We use $v_p(p)=1$.
A nonzero cyclotomic integer satisfies $C(z)\ge1$. Consequently
$C(z)\ge|\Re\langle z\rangle|$ by Cauchy--Schwarz.

\begin{lemma}\label{result:coupled-gap}
Let $A$ be a sum of an even number of roots of two-power order. If $A\ne0$, then $C(A)\ge2$.
\begin{enumerate}[label=\textup{(\roman*)},nosep]
\item If $B$ is a sum of at most $k$ roots, $u>k/2$ is rational,
and $(A,B)\ne(0,0)$, then $C(uA+B)\ge\min\{1,2(u-k/2)^2\}$.
\item Suppose the roots in $B$ have orders supported on $\{2,3\}$.
If $B$ has at most eight terms and $9A+4B$ is a rational linear
combination of roots with total absolute coefficient less than two,
then $A=B=0$.
\item If $B$ as in \textup{(ii)} has at most five terms, $E$ has zero
or two dyadic terms, and the same coefficient bound for $36A+16B+9E$
is less than six, then $(A,B,E)$ is zero or $(0,-\eta,2\eta)$ for
a dyadic root $\eta$.
\end{enumerate}
\end{lemma}
\begin{proof}
Normalized trace onto a dyadic field containing all dyadic parts
multiplies each dyadic part by the rational mean of its odd part.
In its orthonormal power basis the projected $B$ has coefficient
$\ell^1$-norm at most $k$, while the integral vector $a$ of $A$ has
even coordinate sum, so $C(A)=\|a\|^2\ge2$ for $a\ne0$. On $m\ge2$ nonzero coordinates of $a$ the
squared norm of the projected $uA+B$ is at least
$(um-k)^2/m\ge2(u-k/2)^2$. For $m=1$ its bound is $(2u-k)^2$.
For $m=0$ use $C(B)\ge1$.

In (ii)--(iii) write $Z$ for the displayed combination. If $B$ is
outside the dyadic field, its squared distance from it is at least
$3/4$. Indeed an orthogonal component above the cubic-root extension
has norm at least one, while $P+Q\zeta_3$ with dyadic $P,Q$, $Q\ne0$,
has squared distance $3C(Q)/4$ \cite[Lemma~5.2.1]{CMSCyclotomic}.
This gives $C(Z)\ge12$ or $192$, respectively, exceeding the squared
coefficient bound. Thus $B,Z$ have integer coordinate vectors $b,z$
in the dyadic basis. In (ii) $\|z\|_1<2$ and even coordinate sum
give $z=0$. Now $B=-9\delta$, $A=4\delta$ and all conjugates
of the integer $\delta$ have modulus at most $8/9$.

In (iii) $\|b\|_1\le5$ and $\|z\|_1<6$. At $e_j=z_j=0$ the
relation forces $a_j=b_j=0$. For $e=0$, a nonzero $z$ has one
entry $\pm4$, forcing the sole nonzero $a_j$ odd. If $e$ has
one entry, change its sign to $2$. There $z_j=\pm2$ and all other
entries of $z$ vanish. Thus $a_j$ is even and $9a_j+4b_j=-4$
or $-5$, forcing $a_j=0$, $b_j=-1$. For two entries $\pm1$ of
$e$, the absolute $z$-values there are $(1,1)$ or $(1,3)$ up to order.
Modulo $36$, $16b_j=z_j-9e_j$ forces $|b_j|\ge4$ for $|z_j|=1$
and $|b_j|\ge3$ for $|z_j|=3$, exceeding five in total.
\end{proof}

\paragraph{Principal blocks.}
At a place above $p\mid D$, labels form a $p$-block when their
characters have the same reduction on $R$. The partition is independent
of the place by covariance on columns. Over the completed valuation
ring $\mathcal O$ of a splitting field the block idempotent
$e_B=\sum_{j\in B}e_j$ is integral by idempotent lifting
\cite[Tag 04GH]{Stacks}. The character transform identifies $e_j$ with the $j$th coordinate vector, also denoted $v_j$.
Consequently
\begin{equation}
 \sum_{j\in B} d_j^2\chi_j(y)=D\tau(e_By)\in D\mathcal O
 \qquad(y\in R\otimes\mathcal O).
\end{equation}
We use $\chi_j(y)$ also for the linear extension of $\chi_j$.

\begin{lemma}\label{block:height}
The principal block $P=P_p$ contains the vacuum and is Galois-stable.
One has
\begin{equation}\label{eq:principal-trace}
 p^{v_p(f_i)}\mid\sum_{j\in P}|\chi_i(j)|^2,\qquad
 p^{v_p(D/d_i)}\mid\sum_{j\in P}d_j\chi_i(j)
\end{equation}
in $\ZZ$. A row of $p$-power codegree is supported in $P$, with
support equal to $P$ if $p\nmid d_i$. If $p^2\mid D$, then
$p^2\mid f_i$ for every $i\in P$.
\end{lemma}
\begin{proof}
The norm divisibility and height bound are the principal-block
theorems \cite[Theorem~2.11 and Corollary~4.1]{BDPcore} for the
symmetric order $R$, whose unit-coefficient pairing makes the fusion
basis orthonormal against its dual. The norm contains the vacuum
contribution one, so for $p$-power $f_i$ it exhausts the row norm.
If $p\nmid d_i$, then for $j\in P$ the entry $\chi_j(i)\equiv d_i$
is nonzero. The second sum equals $(D/d_i)\tau(e_Px_i)$ by symmetry
and is a rational integer by Galois invariance.
\end{proof}

\section{Blocks and primary phases}
The Chinese remainder theorem factors $\mathrm{SL}_2(\ZZ/L)$ as the product of its prime-power groups. For $p\mid D$ this gives local Weyl, unipotent
and torus operators $W_p,U_p,H_p(a)$, the last corresponding to
$\operatorname{diag}(a,a^{-1})$. The $p$-phase of $i$ is the
$p$-part $\theta_{i,p}$ of its twist. We say that row $i$ meets label $j$ if $S_{ij}\ne0$. A label is active at $p$ if
this phase differs from one. A phase fiber fixes any specified
primary phases. Its mass is the sum of the $d_i^2$ in it. Its
projector commutes with every unspecified local factor.
For a multiplier $a$ supported at $p$, covariance gives
$H_p(a)s=\sigma_a(s)$. The torus acts by the Galois permutation
(up to the common sign $\sqrt D/\sigma_a(\sqrt D)$), and raises
normalized phases to square-unit powers.

\begin{lemma}[Local block rules]\label{block:admissible}
Every full $p$-block spans a local invariant subspace. A row meets
a $p$-active label precisely when $p\mid f_i$. Its coordinate line
is $W_p$-invariant precisely when $p\nmid f_i$. Activity of $i$
requires $p\mid f_i$, and for odd $p$ its $p$-phase order divides
$f_i$. A label of $p$-power codegree belongs to the principal block
$P_p$. Then every $p$-free degree belongs to $P_p$ and has
$p$-power twist. A local invariant coordinate plane on which the
Weyl moves one coordinate line occurs only at two, with equal degrees
and phase ratio of order four.
\end{lemma}
\begin{proof}
For $A_j=sT^js^{-1}$, substituting the primitive idempotents in
\eqref{alg:rotation} gives
\begin{equation}\label{eq:rotation-entry}
 L_j^i(e_z)=|(A_j)_{iz}|^2,\qquad
 (L_j^i(y))^p\equiv L_{j/p}^i(y^p)\pmod{\mathfrak p}\quad(p\mid j),
\end{equation}
the latter for integral coefficients by the rotation certificates.
Here and below we use the ordinary-table arguments of
\cite[\S\S3,6,19]{BDPcore}. To check their hypotheses,
Lemma~\ref{alg:rotation-arithmetic} realizes every divisor trace on
a nonnegative integral element as a power trace of a single multiset
of roots, with all power traces in $\QQ(\zeta_N)$. This is precisely
the rotation input in \cite[Lemma~3.1]{BDPcore}: Frobenius gives the
congruence above, and Galois invariance of multiplicities gives its
higher-power congruence (3.1). Its Theorem~3.2 uses only the diagonal
insertions $i,i^*$, so applies to our $L^i_j$ and gives block
orthogonality. The local action, line and order criteria and the plane
rule follow by the proofs of its Lemmas~3.3--3.4, Theorem~3.6 and
Lemma~19.5. Their remaining inputs are congruence covariance, conductor
divisibility, $\operatorname{rad}N=\operatorname{rad}D$ and the
positive rational vacuum ratios. For prime-power codegree the
complementary Weyls preserve the line. Its $s$-column therefore
lies in its $p$-block. The vacuum entry makes the block principal.
All $p$-free degrees meet this row by reduction, so belong to its
block and have $p$-primary twist by the row criterion.
\end{proof}

\begin{lemma}[Principal support]\label{block:witness}
The principal block $P=P_p$ contains at least three distinct nonunit
degrees. They share a prime in pairs, have gcd one, and every one
meets each $p$-free nonunit degree in a common prime. For any
partition of the primes of $M$ into two or more nonempty parts,
at least three distinct nonunit denominators meet two parts, and
their lcm is $M$.

The $p$-free nonunits of $P$ occupy at least two Galois orbits.
\end{lemma}
\begin{proof}
Orthogonality of an active row to the vacuum on $P$ gives
$\sum_{0\ne Y\in P}d_Y\chi_X(Y)=-1$. The degree gcd follows.
A principal character is nonzero on every $p$-free degree by
reduction. Coprime vanishing then gives the asserted pairwise
intersections and hence at least three distinct degrees. On $P$,
$p\mid b_Y$ by Lemma~\ref{block:height}. Given a prime $t\mid M$,
choose $p$ in another part. Some nonunit in $P_p$ has $t$-free
degree, proving the lcm assertion. If only two crossing degrees
occur they are coprime. For $p$ in one part $J$ all other nonunit
degrees in $P_p$ are multiples of the part of $M$ outside $J$.
Thus exactly one crossing degree belongs to $P_p$, and it is
supported on $J$. Repeating in another part gives a coprime
$p$-free degree, contradicting principality.

For every row $X$, the principal trace
$\sum_{Y\in P}d_Y\chi_X(Y)=D\tau(e_Px_X)/d_X$ is zero modulo $p$.
Its vacuum term is one, so $X$ meets a $p$-free nonunit of $P$.
Burnside vanishing supplies a row vanishing on any one chosen
nonunit orbit, proving the two-orbit assertion.
\end{proof}

\begin{lemma}[Counting phases]\label{block:packets}
Fix an odd prime $p$. In a complementary phase fiber, also after
restriction to a full $p$-block, primitive $p$th roots have equal
mass when $v_p(D)$ is even. For odd valuation the mass at order $p$ is
$(p-1)$ times the neutral mass. At fixed degree and complementary
phases an order-$p^e$ orbit has $p^{e-1}h_p$ distinct phases, where
$h_p=(p-1)/2$.
Put $A=\{t\mid q:t\text{ odd prime}\}$,
$C=\prod_{t\in A}(h_t+1)$, and
\[
 B(q)= C\max\bigl(\{2^{\mathbf1_{2\mid q}}\}
               \cup\{(h_t+2)/(h_t+1):t\in A\}\bigr).
\]
Every nonempty fiber fixing odd-primary phases outside $A$ has
at least $B(q)$ labels. For a set $J$ of odd primes of positive even
valuation, at least $\sum_{p\in J}(p-1)$ nontrivial $J$-phase tuples
occur. At every prime $p\mid D$, including two, the $p$-active mass $V_p$ satisfies
$2pxV_p\ge(p-1)D$ for every active degree $x$ and for the degree
of every based generator.
\end{lemma}
\begin{proof}
Apply the ordinary-table slice formulas of
\cite[Propositions~3.9, 19.1, 19.8 and \S19.2]{BDPcore}. Their
inputs are block invariance and the Bruhat identity on a commuting
fiber projector. The torus supplies the indicated square orbits. For
the last estimate the integer $L^0_{N_{p'}}(X)$ is strictly less
than $x$ if a fusion edge changes $p$-phase. This holds at the unit
for an active label and by connectedness for a generator. Averaging
the roots with weights $N_{ab}^X d_a d_b/D$ bounds the deficit by
$2pxV_p/((p-1)D)$, as in \cite[Proposition~3.8]{BDPcore}.
\end{proof}
A packet consists of $h_p$ equal-degree labels on a square orbit.
Degree $d$ of multiplicity $m_d$ supplies at most $n_d=\lfloor m_d/h_p\rfloor$
packets. At even valuation the differences $k_d$ of the two primitive
packet counts satisfy $|k_d|\le n_d$ and $\sum k_dd^2=0$. If
$n_d-|k_d|$ is odd and either $<p$ or $p^2\nmid f_d$, at least
one place is unused: further primitive packets come in pairs and
a higher orbit costs $p$ places. Two packets alone therefore have
equal degree. At $2\mid q$ the plane rule gives
$\sum_d(n_d+\lfloor n_d/2\rfloor)\ge6$
\cite[Lemma~19.6]{BDPcore}.

Call a prime of $M$ outside $q$ allowed.
\begin{lemma}[Local tensor factors]
Assume $q=1$. For a set $P$ of primes of $M$, let $F$ run over
the maximal subsets occurring together in a twist. For each
$K\subseteq F$ one may simultaneously reserve $\prod_{p\in K}(p-1)$
labels active within $P$ exactly on $K$ whose denominator is
divisible by $\prod_{p\in F}p$. If a label active on $J\subset P$
has $\ell\in P\setminus J$ dividing its denominator, a twist is
active throughout $J\cup\{\ell\}$.
\end{lemma}
\begin{proof}
Decompose the vacuum cyclic module into irreducible tensor products.
Each factor contains a vector $v$ with the vacuum Borel character.
At an allowed odd $p$ this character is trivial. An eigenvalue of order at least $p^2$
has at least $p(p-1)/2$ square-unit translates.
Otherwise the eigenvalues are $p$th roots, and Bruhat makes
$\langle U_p^aW_pv,W_pv\rangle$ constant for $a=1,\ldots,p-1$.
Fourier inversion gives equal weights at all nontrivial roots.
If these vanish, $v,W_pv$ span a $W_p$-stable Borel-fixed space.
Both generating unipotents act trivially. Thus an active factor has
at least $p-1$ active dimensions besides its neutral vector. At two
one active and one neutral dimension suffice.

Every phase occurs in the cyclic module since $sv_0$ has positive
entries. Local phases combine independently in a tensor product,
so distinct maximal sets $F$ require orthogonal summands. Each
summand witnessing $F$ is supported where $\prod_{p\in F}p\mid b_i$,
by orthogonality to the invariant coordinate lines with $p\nmid b_i$.
Its simultaneous active and neutral spaces have the asserted product
dimensions. Their unions of coordinate supports dominate the
corresponding sums of dimensions, giving the reservations by Hall's
theorem. Finally a fixed $J$-phase space without $\ell$-activity
has scalar upper and lower $\ell$-unipotents, hence scalar Weyl.
The line criterion forbids $\ell\mid b_i$ there.
\end{proof}
\paragraph{Integral Weyl columns.}
For odd allowed primes $J$, put $F=\prod_{t\mid M,\,t\nmid q,\,t\notin J}t$
(the product is over primes) and $z=(\prod_{p\in J}W_p)v_0$.
Write $z'=W_J^{-1}sv_0$ for the complementary Weyl column.
Under Galois conjugation each local Weyl acquires a torus factor.
On $J$ the vacuum character is trivial and on the complement it
is the character of $\sqrt q$. Thus $z,\sqrt q\,z'$ are rational.
For either column $x$ the selecting CRT exponent $k$ gives
$f_i|x_i|^2=L_k^0(f_ie_i)$ by Bruhat. This is integral and is
divisible above $t$ for each complementary $t\mid f_i$, by
\eqref{eq:rotation-entry} and $(f_ie_i)^t=f_i^{t-1}(f_ie_i)$.
Since $q$ is square-free, $a_i=b_i z_i$ and $a'_i=\sqrt q\,b_i z'_i$
are integers. The complementary Weyl takes $z$ to $sv_0$ while
preserving each $J$-phase fiber $E$, whence
\begin{equation}\label{eq:phase-column}
 (b_i,F)\mid a_i,\qquad
 \sum_{i\in E}d_i^2=q\sum_{i\in E}d_i^2a_i^2 .
\end{equation}
In each fiber, divide by the squared degree gcd. If some $b_i$ is
divisible by $t\mid F$, the terms at maximal $v_t(b_i)$ sum to
zero modulo $t^2$. They number at least four for $t=2$, three
for $t\equiv3\pmod4$, and two otherwise.

\begin{lemma}[Moments on a fiber]
Let $p$ be odd and allowed. In any complementary fiber, write $J$
for the $p$-active mass in weights $w_i=b_i^{-2}$ and $C_0$ for
the weight at one primitive $p$th root. On the neutral labels with
$p\mid b_i$ there are $a_i\in p\ZZ$ satisfying
\[
 \sum_i(1-a_i)w_i=C_0,\qquad\sum_i(a_i^2-1)w_i=J.
\]
\end{lemma}
\begin{proof}
Take $J=\{p\}$ above and use the complementary coefficients $a'_i$,
writing them as $a_i$. They vanish at active labels and are divisible
by $p$ where $p\mid b_i$. On singleton blocks $p\nmid b_i$ the
local Bruhat relation gives $a_i=1$. Set $v=sv_0$, $z=W_p^{-1}v$.
For the complementary projector $E$, the norm equality
$\|Ez\|=\|Ev\|$ gives the second identity. The vacuum Bruhat
identity for $s^{-1}Es$ identifies $\langle z,Ev\rangle$
with the Galois mean of $\sum_{i\in E}w_i\theta_{i,p}/q$.
The root means $1,-1/(p-1),0$ at orders $1,p,>p$ and equality
of primitive masses give the first identity.
\end{proof}
\begin{corollary}\label{result:mixed}
If two allowed primes $p,t$ occur together, with $p$ odd, at least
$h(t)(p-1)$ active denominators are multiples of $pt$, where $h(2)=4$
and $h(t)=3$ or $2$ as $t\equiv3$ or $1\pmod4$. 
\end{corollary}
\begin{proof}
In a nontrivial $t$-phase, equal primitive masses or the orbit of a
higher phase give at least $p-1$ nontrivial $p$-phases. In each
apply \eqref{eq:phase-column} with $J=\{p\}$ at its maximal
$t$-valuation.
\end{proof}
\begin{corollary}\label{result:fiber-squares}
In square dimension simultaneous activity at three and at $p=5$ or
$7$ requires at least $18$ or $22$ labels, respectively, active at
one of these primes and with $3p\mid b_i$.
\end{corollary}
\begin{proof}
Adding twice the first fiber moment to the second expresses $J+2C_0$
as $\sum((1-a_i)/b_i)^2$ over the eligible neutral labels.
In each complementary fiber meeting $p$-activity there are at least $p-1+s_p$ eligible
labels, where $s_5=3$, $s_7=2$. Indeed a higher phase costs
$p(p-1)/2$ labels. Otherwise each primitive phase has positive equal
mass $C_0$. If each contains at least two labels the required count
already holds. A singleton gives $C_0=b^{-2}$, and the neutral
terms then express $p+1$ as a sum of rational squares. Eight is
not a square, and six is not a sum of two rational squares
(reduce a primitive integral representation modulo three).
At least two nontrivial three-phases meet $p$-activity by equal
primitive masses or the higher-phase orbit. Each gives $p-1+s_p$
labels divisible by $3p$. Conversely at least $p-1$ nontrivial
$p$-phases meet three-activity, and the three-fiber norm in each
requires a further three-neutral label with $3\mid b_i$. This gives
$2(p-1+s_p)+(p-1)$. At equality for $p=7$ the two nontrivial
three-phases each contain six singleton primitive seven-phases and
two seven-neutral labels. The other six are seven-active and three-neutral.
\end{proof}
In square dimension a $p$-phase fiber of at most two
members has $|a_i|=1$ throughout by \eqref{eq:phase-column}.
A vanishing coefficient on a two-label fiber would make $a^2-1$
a nonzero rational square for an integer $|a|\ge2$. Thus without
a $p$-power codegree there are at least $3p$ labels. If $W$ denotes
the mass at each primitive phase and $H$ the mass at higher phases,
$G=\sum d_i^2\theta_{i,p}=D-pW-H$, by square-unit sums.
The vacuum Bruhat identity $|G|=D|z_0|$ gives $MF\mid G$.

\begin{lemma}[Small blocks]\label{block:defectone}
A full $p$-block containing a label with $v_p(f)=1$ has two
defect-one tame inertia orbits with equal complementary phases.
At two it is an equal-degree pair. At odd $p$, let the sizes
and degrees of the neutral and active orbits be $(a,u),(b,v)$.
Then $a,b\mid p-1$, $(p-1)/2\mid b$, and $bv^2=(p-1)au^2$.
At three the blocks have denominators $(z,z,z)$ or $(z,2z,2z)$.
At five $(a,b,v/u)$ is $(1,4,1),(2,2,2)$ or $(4,4,2)$.
\end{lemma}
\begin{proof}
At odd $p$ the defect-one symmetric order theorem and the slice
mass identity give two tame branches, one neutral and one active,
with $bv^2=(p-1)au^2$ also on every complementary fiber
\cite[Proposition~2.13, Corollary~3.10, Proposition~19.11]{BDPcore}.
Square-unit covariance gives $h_p\mid b$ and the lists follow.
At two the defect-one branch is unramified. A second branch of
degree $e=2^k$ has $v_2(f)=k+1$ and order conductor exponent one
\cite[Lemma~2.12]{BDPcore}. For $k>0$ its field is
$F(\zeta_{2^{k+1}})$ over the unramified base $F$, by conductor
divisibility and degree. Its different exponent $ke$ would give
order conductor exponent $e(k+1)-ke=e>1$. Thus both branches are
singletons and the plane rule applies.
\end{proof}

\section{Small codegrees and isolated primes}
For a nonunit row of degree $n$, codegree $f$, write $c_Y$ for
the Galois mean of $|\chi_X(Y)|^2$. The vacuum modular relation,
row norm and lcm integrality give, respectively,
\begin{equation}\label{st}
 \sqrt f\le n^{-1}\sum_Yd_Y\sqrt{c_Y},\qquad
 \sum_Y c_Y=f,\qquad c_Y\ge(d_Y/(d_Y,n))^2\quad(c_Y>0).
\end{equation}
The nonunit hit degrees have gcd one by vacuum orthogonality.

\begin{lemma}[Small codegrees]\label{result:squarefree-codegree}
No codegree is square-free or equal to four or nine. If $f=p^2$,
the nonunit degrees at entries of nonzero absolute trace are prime
to $p$, divide the $p$-free part of $n$, and have gcd one.
For odd $p$, a row of codegree $p^2$ has twist of order at most $p$.
\end{lemma}
\begin{proof}
If $f$ is square-free, act on its coordinate vector by the local
Weyls for $\ell\mid f$ in increasing order. All complementary
factors preserve the line. At a defect-one step the degree changes
by a ratio $k^{\pm1}$ with $k^2\mid2(\ell-1)$, or $k=1$ at two,
by Lemma~\ref{block:defectone}. All remaining prime valuations are
unaffected. At a path to the vacuum, $n$ therefore divides a product
of these $k$'s. For $f>1$ its largest prime then has valuation one
in $D$, impossible by square-fullness. For $f=1$ this gives $D=1$.

At $f=4$ all hits lie in $P_2$, where nonunit degrees share a prime
in pairs. Their gcd is one, so there are at least three distinct
nonunit hit degrees. The norm budget forces exactly three hits,
each with $c_Y=1$. All three divide $n$ by \eqref{st}. If one
equals $n$, the other two would be coprime. Thus all three are
proper, and $2\le1/n+3/2<2$ by the first inequality in \eqref{st}.

At $f=p^2$ the nonzero-trace degrees have gcd one by the
trace-vector identity \cite[Proposition~5.4]{BDPcore}. Each
nonunit absolute trace $t$ has $|t|<p$. Its square costs at least
$t^2$ in the projected-row identity \cite[(2.4)]{BDPcore} besides
the vacuum. Principality gives $t\equiv0\bmod p$ when $p\mid d_Y$,
since the row field is totally ramified at $p$. Thus all the
nonzero-trace degrees divide the $p$-free part of $n$.

If a row of codegree $p^2$ has twist of order $p^2$ at odd $p$,
its field contains $\QQ(\zeta_{p^2})^+$ by square covariance.
Write its degree as $pm$, $m\ge(p-1)/2$. The wild and full Galois
means $v,w$ of the row have squared norms $p,p/m$ by orthogonality.
Wild averaging a $p$-power
root fixes it or kills it, so each entry $\beta$ of $v$ is
$\sum_{k=1}^{p-1}c_k\zeta_p^k$ with integral coefficients. Then
\[
 \bigl(C(\beta)+(p-1)|\langle\beta\rangle|^2\bigr)/p
       =\sum_k c_k^2/(p-1).
\]
Multiplied by the column orbit size, the right side is integral.
The $c_k$ are constant on stabilizer cosets, whose index divides
that size. At least four orbits contribute positively, since the
nonunit degrees of nonzero mean have gcd one, lie in $P_p$, and
share a prime in pairs. Hence $4\le (\|v\|^2+(p-1)\|w\|^2)/p=1+(p-1)/m\le3$,
a contradiction.

A codegree-nine label would have denominator three, excluded
by Proposition~\ref{result.One-packet} below using only the prime-square
estimates just proved.
\end{proof}

\begin{lemma}[Reduction to even dimension]\label{general:odd-small}
If $r\le44$, the global dimension is even.
\end{lemma}
\begin{proof}
Suppose $r\le44$ and $D$ is odd. Let $b=\min_i b_i$. Since $q=\sum_i b_i^{-2}$,
$qb^2\le44$. The denominators are odd, and $b=1$ is excluded by the
square-free codegree bound (Lemma~\ref{result:squarefree-codegree}).

If $q=1$, then $b=3$ or $5$. In a row of codegree $b^2$, the nonunit
degrees of nonzero trace are prime to $b$ and have gcd one
(Lemma~\ref{result:squarefree-codegree}). Each has at least two prime divisors, so
at least three primes other than $b$ divide $D$. By
Lemma~\ref{block:packets} there are at least
$(3-1)+(5-1)+(7-1)+(11-1)=22$ nontrivial phase tuples.
Each contains a dual pair, giving $r\ge1+2\cdot22$.

If $q>1$, then $b\ge3$ forces $q=3$. Two other primes $p<t$ divide
$D$. Lemma~\ref{block:packets} gives at least $p+t-2\ge10$ nontrivial
$(p,t)$-phase tuples, each occurring on at least $B(3)=3$ labels.
Every such label has denominator divisible by $p$ or $t$, hence at
least five. The mass is therefore at most $14/9+30/25<3$.
\end{proof}

\subsection{Weyl columns in square dimension}
Assume $q=1$. Join two primes when they divide the order of a common twist.
A nontrivial phase on a component of this graph has only
vacuum phases outside that component. The complementary upper and
conjugate lower unipotents act by their vacuum scalars on this
fiber. Their Weyls act by scalars, so the codegree there is
supported on the same component.

\begin{lemma}[Component columns]\label{result:component-column}
Let $J$ be a nonempty union of odd twist components, $W=W_J$, $w=Wv_0$.
Write $M_J$ and $u_i$ for the $J$-parts of $M$ and $b_i$,
respectively. Then $\epsilon=\gamma^3=\pm1$ and $g=1-w_0$ obeys
\[
 g=\sum_{\theta_{i,J}\ne1}b_i^{-2}
     \left(1-\frac{\mu(m_i)}{\varphi(m_i)}\right)>0,\qquad
 (M/M_J)(M_Jg)\equiv-\epsilon\pmod {M_J},
 \quad m_i=\operatorname{ord}(\theta_{i,J}).
\]
Its reduced denominator is $M_J$. At $i\ne0$,
$\epsilon w_i=A_i/u_i$ with $|A_i|<u_i$ and
$A_i(b_i/u_i)\equiv1\pmod{u_i}$. When $b_i$ is supported on $J$,
$\epsilon w_i=1/b_i$. Moreover
\begin{equation}\label{eq:component-norm}
 \sum_{i\ne0}w_i^2=2g-g^2,\qquad
 \sum_{i\ne0}w_i/b_i=g(\epsilon+1/M).
\end{equation}
For $J=\{p\}$ the principal $p$-block is $\{i:p\mid b_i\}$.
\end{lemma}
\begin{proof}
Write $V$ for the complementary Weyl. The split-column identities
\cite[Theorem~6.1, Lemma~6.2, Proposition~6.3]{BDPcore} use only
\eqref{eq:rotation-entry}, congruence covariance and Bruhat. In
particular $s=V+\epsilon(W-I)$, $W^{-1}v_0=w$, and the rational
columns $w,Vv_0$ are integral outside their respective prime sets.
On an active coordinate $Vv_0=0$, so $\epsilon=\pm1$.
The vacuum Bruhat sum, averaged over units, gives the displayed
Ramanujan formula. Comparing the split column with
$(sv_0)_i=1/b_i$ gives the congruences, including the exact
denominator at the vacuum (also $b_iw_i\in\ZZ$ by
\eqref{eq:phase-column}). Thus $w_0\ne0$ and $|w_i|<1$. A pure
coordinate line is invariant under $V$, so $|w_i|=1/b_i$ and
the congruence fixes its sign. The nonzero entries and the
principal height bound now give the singleton block assertion.
Finally $\|w\|=1$ and $\langle(W-I)v_0,sv_0\rangle=\epsilon g$
give \eqref{eq:component-norm}.
\end{proof}

\begin{corollary}\label{result:mixed-neutral-rank}
With at least two twist components, at least three distinct nonunit
degrees occur on twist-one labels whose denominators meet two
components.
\end{corollary}
\begin{proof}
Apply Lemma~\ref{block:witness} to the partition by components.
An active label has its entire codegree supported on its own component.
\end{proof}

At an isolated odd prime $p$, let $E$ be the reciprocal-square
mass of its $h$ active labels, all of denominator at least $B$. The
active entries use $E$ of the squared norm $2g-g^2$, where
$0<g\le pE/(p-1)$. On a mixed denominator with $p$-part $u_i>1$
the entry has magnitude at least $1/u_i$. Thus
\begin{equation}
 \sum_{i\ne0:\ 1<u_i<b_i}u_i^{-2}
       <\frac{p+1}{p-1}\frac{h}{B^2}.
\end{equation}

\begin{lemma}[Sparse activity]\label{result:activity-height-gap}
If $h<p(p-1)$ and $p^a$ is the largest active denominator, then
\begin{equation}\label{eq:sparse-height}v_p(M)=2a-1.\end{equation}
Each nontrivial $p$-part of a nonunit denominator is at least the
least active denominator $B$.
\end{lemma}
\begin{proof}
At the largest level where the multiplicities of the two primitive
square classes differ, their difference is divisible by $p^2$ by
equality of their masses. This costs at least $p^2(p-1)/2>h$ labels.
Thus all levels are balanced. A higher orbit has size at least
$p(p-1)/2$, so at most one occurs, at order $p^2$. At each level
$p^i$ the charge (its contribution to $g$) is $pc_i/p^{2i}$, where
$c_i=h_{i,1}/(p-1)+h_{i,>}/p$, with counts at order $p$ and above.
Balance and orbit sizes give $c_i\in\ZZ$ and $(p-1)c_i\le h$.
In particular $0<c_a<p$, so the highest level is the unique term
of lowest valuation. The denominator of $g$ proves \eqref{eq:sparse-height}.

Write $h=(p-1)j+p(p-1)k/2$, where $k=0$ or $1$. Then
$j\le p-1$ or $(p-1)/2$, respectively, so
$2g-E\le((p+1)j+p(p-1)k/2)/B^2\le(p^2-1)/B^2$.
An inactive denominator with $p$-part $1<u<B$ would use at least
$1/u^2\ge p^2/B^2$ of the remaining norm $2g-g^2-E$.
\end{proof}

\begin{lemma}[Valuation one]\label{result:weyl}
Suppose $p\Vert M$ is odd and isolated, and set $n=M/p$. The
$p$-active labels are of degree $n$. Each primitive phase has the
same multiplicity $0<j<p$. Off the vacuum set $a_i=p\epsilon(w_p)_i$.
Then
\begin{equation}
 a_i\equiv d_i/n,\quad nj\equiv-\epsilon\pmod p,\qquad
 \sum a_i^2=2pj-j^2,\quad \sum (d_i/n)a_i=pj+\epsilon j/n .
\end{equation}
The nonzero coefficients occur exactly at $p$-free degrees. Their
proper nonunit degrees have gcd one and share a prime in pairs.
\end{lemma}
\begin{proof}
Isolation gives active codegree $p^2$, hence order $p$ by the tame
row bound in Lemma~\ref{result:squarefree-codegree}. The primitive masses agree,
so $g=j/p$ and the component lemma gives the congruences and sums.
The trace support of a prime-square row has gcd one and consists
of divisors of $n$, so its proper degrees still have gcd one and
are pairwise noncoprime by principality. The active mass gives
$j\le p+1$. At $j=p+1$ the $(p-1)j$ active unit coefficients
exhaust $2pj-j^2$, leaving no proper hit. Finally $p\nmid j$ by
the denominator of $g$.
\end{proof}

\begin{proposition}[A prime denominator]\label{result.One-packet}
If a prime denominator $p$ occurs then $p\ge5$, $p\Vert M$ and
$p$ is isolated. In the notation of Lemma~\ref{result:weyl}, $j$ is
a proper divisor of $p-1$ and exactly $pj$ labels have degree $n$.
Their $j$ neutral members form a single Galois orbit. At every
proper nonunit $p$-free degree let $k_Y=n/d_Y$. Then
\[
 \sum a_Y^2=j(p-j),\quad \sum a_Y/k_Y=\epsilon j/n,\quad
 k_Ya_Y\equiv1\pmod p,\quad \lcm_Y k_Y=n .
\]
For
each fixed primitive phase, with label set $A$, and each neutral nonunit $Y$,
\begin{equation}\label{eq:packet-trace}
 \sum_{X\in A}\chi_X(Y)=-\epsilon a_Y.
\end{equation}
For a neutral $p$-inertia orbit of length $h\le(p-1)/2$ one has
$j\mid h a_Y$.
\end{proposition}
\begin{proof}
Codegree four is excluded, so $p$ is odd. The height bound
(Lemma~\ref{result:height}) gives $p\Vert M$. Every active label
has $p$-free degree, whence $p$-primary twist by the prime-power
row rule. Thus Lemma~\ref{result:weyl} applies, and the local level
is $p$ since all active phases have order $p$.
All $m$ coefficients of degree $n$ equal one by the pure-denominator
identity of Lemma~\ref{result:component-column}.
Put $H=j(p-j)$ and $z=pj-m$. After removing this class,
\begin{equation}\label{eq:weyl-cost}
 \sum a_Y^2=H+z,\qquad \sum a_Y/k_Y=z+\epsilon j/n,qquad
 -\frac{H+z}{p-1}\le z+\epsilon j/n\le\frac{H+z}{p+1}.
\end{equation}
The bounds use $k_Ya_Y\equiv1\bmod p$ and $k_Y>1$.

Write $Q$ for the neutral projector and $u=\sum_{X\in A}v_X$.
The phase projection of $W_pv_0$ is $\epsilon u/p$. Thus $u$
lies in the Steinberg summand of the local vacuum cyclic module,
a quotient of the permutation module on $\mathbb P^1(\mathbb F_p)$,
whose nonconstant part is irreducible by double transitivity. The Steinberg neutral line
is spanned by $g=Q(W_p-I)v_0$, with $g_0=-j/p$. The complementary
Weyl is scalar $\epsilon$ on $u$, so $Qsu$ lies on this line.
Its vacuum entry $j/p$ gives $Qsu=-g$, proving \eqref{eq:packet-trace}.

For an inertia orbit $O$ of neutral labels of degree $d$ and
size $h\mid p-1$, put $b_X=\sum_{Y\in O}\chi_X(Y)$. These are
integers since the active row fields lie in $\QQ(\zeta_p)$, and
they are all $hd\bmod p$ by principality. In each packet their
sum is $-\epsilon h a_Y$. If $j\nmid h a_Y$, at least $j-1$
pairs differ by at least $p$, so the squared sum per packet is at
least $((h a_Y)^2+(j-1)p^2)/j$. Over all packets unitarity leaves
at most $hp^2-h^2(d/n)^2$. Thus $j\mid h a_Y$ for $h\le(p-1)/2$.
If $d=n$, $a_Y=1$, $h=p-1$, the opposite assumption would give
$(p-1)^2+(j-1)p^2\le j(p^2-p+1)$, hence $j\le2$, but both possible values divide
$p-1$. Consequently every such $h$ is divisible by $j$, and $j\mid m$.

By \eqref{eq:weyl-cost}, $|z|<2j$. At $z=\eta j$, $\eta=\pm1$,
its signed bound forces $\epsilon=-\eta$ and $jn\le p+\eta$,
hence equality by $jn\equiv\eta\bmod p$ and $n>2$.
Equality in the bound gives $a_Y=-\epsilon j d_Y$, contrary to
\eqref{eq:packet-trace} and $|\chi_X(Y)|<d_Y$. Thus $m=pj$.
The $j$ neutral members form one inertia orbit, so $j\mid p-1$.
If $j=p-1$, a proper orbit $h<j$ costs
$ha_Y^2\ge j^2/h>j$ by divisibility, while $h=j$ alone exhausts
the budget $H=j$. At least three proper degrees occur, so $j<p-1$.
This also excludes $p=3$, where $j=1$ would leave budget two.
\end{proof}

\begin{lemma}[Pure denominators]
An isolated odd $p$ has at least $p$ nonunit labels of pure $p$-power
denominator. If $p$ itself is absent and exactly $p-1$ labels are
active, their common denominator $B$ satisfies $M_p=B^2/p$.
\end{lemma}
\begin{proof}
If $p$ occurs the prime-class formula suffices. Otherwise there are
at least $p-1$ active labels by the primitive masses or the orbit of
a higher phase. At equality each primitive phase is a singleton of
common denominator $B\ge p^2$. Thus $g=p/B^2$ and $M_p=B^2/p$.
On a mixed denominator $uk$ with $p$-part $u>1$ the signed numerator
$A$ satisfies $|A|k\ge u-1$, so $|w_i|/b_i\le w_i^2/(B-1)$ by
the sparse bound $u\ge B$. Subtract the pure contributions from both
identities in \eqref{eq:component-norm}. If $H$ is the inactive pure
mass multiplied by $B^2$, this gives
\begin{equation}\label{eq:packet-shell}
 |1-H+\epsilon p/M|\le\bigl(p+1-H-(p/B)^2\bigr)/(B-1).
\end{equation}
Thus $H\ge1-1/(p-1)-1/(6p)>1/p^2$, since $M\ge6B$:
there is an additional pure nonunit.

More generally, if $h<p(p-1)/2$ then all phases are primitive and
every mixed nontrivial $p$-part is at least $p^2$. For active mass
$V$ and inactive pure mass $P$, the same subtraction gives
\[
 P\ge g-V-g/M-(2g-V)/(p^2-1)
   =\frac{V}{p-1}\left(1-\frac pM-\frac1{p-1}\right)>0.
\]
\end{proof}

\section{Multiplicity and prime support}
Equality in $\langle |z|^2\rangle\ge1$ for a nonzero cyclotomic integer
holds precisely for roots of unity. For a row principal at two, if all its entries on an odd-degree
class of size $s$ have modulus one, their sum is an integer of absolute value
at most $s$ and parity $s$. It is therefore a sum of $s$ signs. If it
equals $\pm s$, every entry equals that sign.

\begin{lemma}[A two-column Gram bound]\label{lem:congruent-rows}
Let a Galois-stable family of $m$ orthogonal rows of cyclotomic integers have norm
squared $f=p^a$ and vacuum entry one, and suppose they are perpendicular to the degree
row. Suppose they vanish off a common coordinate set $F$ and are
congruent to the degree row above $p$ on $F$. Then
$m\le f-\lceil\sqrt p\rceil$ and, for $a\ge2$, $m\ne f-p$.
\end{lemma}
\begin{proof}
Some $p$-free nonunit $Y\in F$ occurs, by degree orthogonality modulo $p$.
Let $w,A$ be the sums of the entries and their squared moduli there.
Both are integers and $A\ge m$ by the norm inequality. Set $k=f-m$,
$l=f-A$. The Gram matrix of $e_0,e_Y$ after projection off the rows,
multiplied by $f$, has diagonal $k,l$ and off-diagonal $-w$.
It is positive definite because proportionality would give a norm ratio
$\sqrt{l/k}\le1$, whereas the ratio of the degree components is $d_Y>1$.
Since $w\equiv-kd_Y$, $l\equiv kd_Y^2\pmod p$, one obtains
\begin{equation}\label{eq:saturation-gram}0<p\mid kl-w^2\le k^2.\end{equation}
If $k=p$, then $l>0$, $p\mid l$ and $l\le p$, whence $l=p$, $w=0$.
This holds at every $p$-free nonunit. At every other coordinate the
row sum is divisible by $p$, contradicting degree orthogonality modulo
$p^2$, since the vacuum sum is $p^a-p$.
\end{proof}

\begin{lemma}\label{result:saturation}
At degree $n$ of prime-power codegree $f=p^a$, the preceding multiplicity
bounds apply. For odd $n$ and dyadic $f$ one has $m\le f-6$ at $f=8$
and $m\le f-8$ at $f\ge16$. For odd $n$ at $f=16$, either $\omega(n)=3$ or $r<38$ implies $m\le6$.
\end{lemma}
\begin{proof}
Use Lemma~\ref{block:admissible} for the first assertion. Suppose
$f-m\le5$ at eight or $f-m\le7$ at $f\ge16$. Every even hit costs
at least $4m\ge f$ in its column, by lcm integrality, besides the
vacuum. Thus $P_2$ is exactly the odd-degree set, all of whose degrees
divide $n$. Its mass is integral at two, so $8\mid |P_2|$, since
odd squares are one modulo eight. The proper nonunit hit degrees have
gcd one and are pairwise noncoprime. Hence at least three distinct
ones occur. In particular $f-8<|P_2|\le f$. Thus $|P_2|=f$ and
every entry is a root of unity.

Put $h=f-1-m$. Degree orthogonality for each row takes the form
$1+nt+\sum_{i=1}^h\epsilon_i a_i=0$ with signs $\epsilon_i$,
$t\equiv m\pmod2$, $3\le a_i\le n/3$ the proper degrees, $n\ge15$.
For $h=4$, $t$ is odd and mixed signs give absolute proper sum
at most $n-3$, so all four signs agree and their sum is determined.
For $h=3,5$, one has $t=0$. Two solutions differ in an even number
of positions, and four differences would leave a signed degree $-1$
(or be impossible at $h=3$). Two differences only exchange equal
degrees. Thus the signed sum $t_b$ at each proper degree $b$ of
multiplicity $s_b$ is constant across rows, with $\sum|t_b|\ge3$.
Bessel on the vacuum and class-constant lines yields
$m(1+9/h)\le m(1+\sum t_b^2/s_b)\le f$, impossible in the asserted
ranges. For $h=6$, $t$ is odd. Bessel on the degree vector restricted
to the vacuum and proper coordinates gives
$m/f\le n^{-2}(1+\sum a_i^2)\le1/225+4/9+2/25<9/16$ if at most
four of them equal $n/3$. Otherwise five equal $a=n/3$, and the
remaining degree, coprime to $a$, must be three. Its singleton entry
is a sign, forcing $a\mid1\pm3$, impossible.

For the final bound assume $m\ge7$. Again $4m>16$ gives the same
sixteen odd labels. Four odd primes would reserve at least
$2+4+6+10$ more labels, so precisely three occur under either
hypothesis. The proper labels omitting each of these primes form
three disjoint nonempty sets. On a set omitting $p$, its signed
degree sum is $-1\pmod p$. On one or two labels this fixes the
signs, since changing both negates the sum. Actual entries are
fixed as well. Two fixed columns would give $3m\le16$. With at
most eight proper positions the remaining possibility includes a
singleton and a triple. On a triple two solutions differ in exactly
two positions, exchanging equal residues after orienting them up to
sign. Thus each oriented residue class has fixed sum. At least one
is nonzero. Projection onto the vacuum, singleton and these class
lines gives $m(2+1/3)\le16$.
\end{proof}

\begin{proposition}[Rows obtained from a defect-one block]\label{result:effective}
In an equal-degree defect-one $\ell$-block $B$ of codegree $\ell p^a$
($p\ne\ell$), there is a unique $\ell$-neutral label $X$. Its character
restricted to $F=\{Y:\theta_{Y,\ell}=1\}$, extended by zero, equals
$v_B=\ell^{-1}\sum_{i\in B}\chi_i$. These rows have squared norm $p^a$,
are pairwise orthogonal, supported on $P_p$, and have character field
$K_X\subseteq\QQ(\zeta_{p^a})$. Lemma~\ref{lem:congruent-rows}
applies to them. Their vacuum Gauss identity is
\begin{equation}
\sum_Yd_Y\theta_Yv_B(Y)=\eta_B\sqrt{D/\ell}\,\theta_X^{-1},
\qquad |\eta_B|=1 .
\end{equation}
The factor is constant on dual pairs and under odd $p$-inertia.
At squared norm $p^2$, the nonunit degrees in the nonzero-mean
support of an actual or such an averaged row are $p$-free, have gcd
one and are pairwise noncoprime. At least three distinct degrees
and at least three primes other than $p$ occur.
\end{proposition}
\begin{proof}
Write $s=W_\ell W_pV$ and $c=\ell^{-1/2}\sum_Be_i$.
Lemma~\ref{block:defectone} and Lemma~\ref{block:admissible} give uniqueness
of $X$, with complementary phases constant on $B$. The projection
of $se_0$ on $B$ spans $c$. Its inverse $W_\ell$-image is neutral,
so $W_\ell^{-1}c$ spans $e_X$. The central square $W_\ell^2$ preserves
that line. The factors of $V$ preserve it too, by
Lemma~\ref{block:admissible}. Hence $X\in P_p$ and $sc$ is supported
on $P_p\cap F$. $s(c^\perp\cap\mathbb C B)$ is supported off $F$.
This gives the stated restriction and the norm by unitarity.
The Galois stabilizers of $X,B,v_B$ coincide, as $s$ separates
the block sums. Inertia at $\ell$ preserves the residue block and
fixes $X$, so the conductor bound removes $\ell$. Average the vacuum
Gauss equations over $B$. The local inverse twist sum is $1\pm i$
at two and a signed quadratic Gauss sum at odd $\ell$, each of modulus
$\sqrt\ell$. The other factors are fixed under odd $p$-inertia.
For the last assertion the trace and congruence argument of
Lemma~\ref{result:squarefree-codegree} applies on $F$ verbatim. Its hit labels
belong to $P_p$, where $\chi_Y(Z)\equiv d_Z\ne0$ above $p$.
Coprime vanishing therefore gives the pairwise property. A pairwise
intersecting family of prime supports with empty intersection requires
at least three primes and three members.
\end{proof}

\begin{lemma}[The height at a prime-square codegree]\label{result:height}
An actual or averaged prime-square row of odd prime $p$ has degree
$n$ prime to $p$. In particular a denominator $p$ in square dimension,
or in an equal $\ell$-block at $D=\ell M^2$, forces $p\Vert M$.
\end{lemma}
\begin{proof}
Sum the rows of its orbit to an integer vector $v$, of squared norm
$hp^2$ and vacuum entry $h$, where $h$ is the orbit size. Only
$p$-inertia is needed. Each nonzero coordinate has $p$-primary twist
and degree dividing the prime-to-$p$ part of $n$ by the preceding
proposition, principal height and local admissibility. On the orbit
the twists are trivial or of order $p$ by Lemma~\ref{result:squarefree-codegree}.
Its norm argument applies also to the averaged translates, whose mean
support has three pairwise noncoprime degrees of gcd one.
The summed Gauss equations have right side
$\eta pn\sum_{X\in O}\theta_X^{-1}$. If $p\mid n$ and the orbit is
neutral, taking absolute values gives
$ pnh\le h+(n/p)(hp^2-h^2)<pnh$, since $n/p>1$.
For an active orbit apply the $p$-Galois average weighted by the
quadratic character. It annihilates roots of order different from
$p$, and sends $\zeta_p^k$ to $(k/p)\tau/(p-1)$ for the quadratic
Gauss sum $\tau$. The inverse row twists belong to one square class.
$\eta$ and $v$ are fixed. Cancelling the common factor gives
$pnh\le (n/p)(hp^2-h^2)<pnh$.
\end{proof}

\begin{lemma}\label{result:small-effective}
At codegree $4\ell$ ($\ell$ odd prime) the multiplicity is one of
$h,2h,\ldots,7h$, $h=(\ell-1)/2$, and there is no equal $\ell$-block.
At codegree eighteen the degree is three-free, its multiplicity at most
ten, and at least four dimension primes occur. In square dimension a
prime denominator $p$ satisfies $p\ge5$, $p\Vert M$ and at least four
primes divide $M$. All $p$-active labels have denominator $p$.
\end{lemma}
\begin{proof}
An equal block at $4\ell$ would yield an averaged norm-four row.
The norm-four argument in Lemma~\ref{result:squarefree-codegree} applies to its principal support and vacuum Gauss identity.
In an unequal block the degree cannot be the smaller. Its square
ratio would divide $4\ell$, giving the forbidden codegree $\ell$.
Thus all these labels are active in packets divisible by $h$. Their
blocks have mass $mn^2\ell/(\ell-1)<4\ell n^2$ since they omit the
vacuum. At eighteen the dyadic blocks are equal pairs, and the
preceding lemma gives $3\nmid n$. Each averaged row meets the neutral
member of each pair and at least three distinct proper nonunit degrees
of gcd one, so $1+m/2+3\le9$. The assertion on the number of primes
follows from Proposition~\ref{result:effective}. In square dimension
also use Lemmas~\ref{result:squarefree-codegree}, \ref{result:squarefree-codegree} and~\ref{result:weyl}.
\end{proof}
Thus for $D=5M^2$ with $9\mid M$, all denominator-three labels are in
unequal quinary blocks, which have as many denominator-six labels.

\begin{lemma}[Exclusion of codegree eight in small rank]
The greatest degree $x$ cannot have codegree eight in rank at most $44$.
\end{lemma}
\begin{proof}
Here $D=8x^2=2M^2$. Its row meets all odd labels, with degrees dividing
$x$. Their number is divisible by four. If $x$ is odd this is the
whole support, with at least three distinct proper divisors, whence
$x\ge105$. There are eight entries, all roots of unity. There are
$m\le2$ maximal labels by Lemma~\ref{result:saturation}. At $m=1$ the
vacuum Gauss bound would give $\sqrt8\le1/x+1+4/3+1/5+1/7$.
At $m=2$ each maximal row field lies in $\QQ(\zeta_8)$ and has
at most two conjugates, so its root entries belong to $\{\pm1,\pm i\}$.
The rows are rational or complex conjugate, hence agree or oppose at
each hit. In each row the two maximal entries have even integral sum,
zero since all other degrees total at most $1+5x/3<2x$.
By symmetry both maximal hits belong to the half-difference.
Row orthogonality gives four hits in each half-row. Degree orthogonality
on the half-difference makes its two proper degrees equal to some $y$.
The half-sum has three proper degrees with signed sum $-1$. A repeated
degree $a$ would force the third to be $1$ or $2a\pm1$, contrary to
pairwise noncoprimality. The two Gauss half-sums are therefore bounded
by $1+(1/3+1/5+1/7)x<7x/10$ and $2x+2y\le8x/3$. Their squared
moduli must sum to $8x^2$, a contradiction.

For even $x$, the norm-eight calculation of
\cite[proof of Theorem~20.7]{BDPcore} uses precisely principality,
$(x/d_Y)\chi_X(Y)\equiv0$ above two and \eqref{st}. It gives
$\chi_X=A+B\sqrt2$, where outside the vacuum $A$ consists of
three signs at the odd nonunits, and $B$ of two opposite signs at
one degree $y$. The odd degrees are distinct, divide $x/2$, have
gcd one and share primes in pairs, so at least three odd primes
divide $x$. In particular $y$ is even. If $y<x$ then
$\sqrt8\,x\le1+(1+1/2+1/3)x/2+\sqrt2\,x$, impossible.
The signs at the odd labels are unique by degree orthogonality
(a change of signs would have to cancel two of the three distinct
odd degrees). Bessel gives $4m\le8$. The conjugate rows give
$m=2$. These two rows and the four odd labels have purely dyadic twist.

Odd-prime activities require $2(p-1)$ disjoint positions per prime
besides the two maximal and four odd labels. Thus the odd primes are $3<l<t$.
Four primes cost at least $44$ places, and $(5,7,11)$ costs $40$.
Each odd nonunit omits a different prime by gcd one and pairwise
noncoprimality, giving odd denominator floors $12,4l,4t,12lt$.

Write $2k$ for the multiplicity of denominator three. If $k>0$ then
$3\Vert M$. The $k$ dyadic pairs yield orthogonal effective norm-nine
rows, all meeting exactly the dyadic-neutral three-free labels, of
degrees dividing $n=M/3$. At $0$ and the unique three-free odd
nonunit $U$ their entries are $1,-A_U$ by orthogonality to $A$,
the mean of the maximal rows. Their Galois-invariant sum is integral,
congruent to $k d_Y$ modulo three at every hit $Y$, and has squared
norm $k(9-2k)$ off $0,U$. Thus $k\le4$. At $k=4$ the four neutral
$n$-labels exhaust this norm, yet any further hit would be nonzero
modulo three. The nonunit hit degrees would then have common divisor
$d_U\mid n$. At $k=3$ every entry is divisible by three, so the
sole remaining hit has degree $|1-A_Ud_U|$ by degree orthogonality,
coprime to $d_U$ in the principal three-block. Thus $k\le2$.

Reserve $2k$ threes and $4-2k$ sixes, and for $s=l,t$ reserve
$2(s-1)$ floors $s$, doubled if denominator $s$ is absent. Fill
the remaining places to $44$ at floor four. The mass bound decreases
with each prime and increases with $k$. For $(l,t)\ne(5,7)$ it is
at most $11977/6600$ (take $5,11$, $k=2$). For $(5,7)$ with at
most one of $3,5,7$ occurring as denominators, its maximum is
$3681/1960$ at $k=0$ and $57959/29400$ at $k>0$. All are below two.

Thus two primes $a,b$ in $\{3,5,7\}$ occur as denominators and
have valuation one in $M$. Writing $t$ for the third, the odd
degrees are $ab,at^i,bt^j$, with $|at^i\pm bt^j|=ab\pm1$.
For $(a,b)=(3,5),(3,7),(5,7)$ the $t$-divisible right side is
$14,20,36$, respectively. If $i\ne j$, dividing by the full $t$-part
leaves $2,4,4$, respectively, each less than $at-b$. Thus $i=j=1$. The signs are $-,-,+$ at $15,21,35$.
Choose $p\in\{5,7\}$ among the denominators. Every degree in $P_p$
is $p$-free by height, and their nonunit gcd is one. Some $Y\in P_p$
therefore omits both five and seven. Its prime support is $\{2,3\}$.
The singleton odd rows are rational. By symmetry, integrality and
coprime vanishing, $\chi_Y(15)=5u$, $\chi_Y(21)=7v$, $\chi_Y(35)=0$,
where $u,v\in\ZZ$ and $|u|,|v|<3$. Orthogonality to the maximal
mean gives $5u+7v=1$, impossible modulo seven.
\end{proof}

\subsection{Prime denominators in square dimension}
Suppose $D=M^2$.
\begin{lemma}
A denominator-four row of even degree or multiplicity at least four
hits only $4\mid b_Y$. Every full prime power in $M$ divides the
denominator of some nonunit hit. With three support primes a denominator
$p^a$ forces $p^2\mid M$, and four fours force $8\mid M$.
\end{lemma}
\begin{proof}
The dyadic column estimate and the height bounds are
\cite[\S21.2 and Lemma~21.6]{BDPcore}. The common-divisor assertion
also follows from gcd one on the hit degrees.
\end{proof}

\begin{lemma}\label{result:five-reduction}
If $v_p(D)>0$ is even, $p$ odd, and all $p$-phases have order at most
$p$, the number of active labels is at least $(p-1)(1+c/2)$, where
$c$ is the number of neutral labels less their number of inertia orbits.
For a prime denominator write $j$ for its multiplicity per primitive
phase. At five, $j=1$ and the four proper Weyl coefficients are signs.
At seven, $j=1$ or $2$. At $j=2$ the proper coefficients have magnitudes
$1,1,2,2$, the first two on a moved pair and the last two on distinct
singleton degree classes.
\end{lemma}
\begin{proof}
For $G=\mathrm{SL}_2(\mathbb F_p)$, its unipotent $U$ and its Borel
$B$, $c=\dim V^U-\dim V^B$ by the permutation action of the torus.
The vacuum Steinberg supplies $p-1$ active dimensions. The other
nontrivial torus types in $V^U$ occur by reciprocity in
$\operatorname{Ind}_B^G\alpha$. A non-self-inverse character gives
two torus lines and $p-1$ active dimensions, and each of the two
quadratic constituents gives one line and $(p-1)/2$ active dimensions
\cite[Prop.~2.14]{PrasadSL}, \cite[Lem.~4.3]{Nevins}.
Indeed the induced module restricts to $U$ as
$\mathbf1\oplus\mathrm{Reg}_U$. Summing proves the bound.

The neutral degree-$M/p$ orbit uses $j-1$ of the at most $2(j-1)$
units of deficit. Each proper orbit of length $h$ and coefficient
$a$ costs $ha^2$ out of $j(p-j)$, with $j\mid ha$
(Proposition~\ref{result.One-packet}). At least three degrees occur.
At $(p,j)=(5,2)$ at most one pair moves, and the norm six pays at
most one other even coefficient. Thus $j=1$, giving four signs.
At $(7,3)$ a pair costs at least eighteen, and a triple and two
fixed labels cost at least twenty-one, against norm twelve. At
$(7,2)$ norm ten on at least three degrees forces the indicated
pair and two even singletons. The possible $j$ are the proper
divisors of $p-1$ by Proposition~\ref{result.One-packet}.
\end{proof}

\begin{lemma}[The least possible activity at an isolated prime]\label{result:packet-idempotent}
Suppose an isolated odd prime $p$ has just $p-1$ active labels,
all at denominator $B$. If $B\ne p$ suppose fewer than
$p^2(1-1/(p-1)-1/(6p))$ inactive nonunit labels have pure $p$-power
denominator. Then the class at $B$ has $p$ labels. At $B=p$ all its
rows take the common value $-\epsilon a_Y$ on each proper $p$-free
label $Y$.
\end{lemma}
\begin{proof}
For $B\ne p$, the remaining Weyl norm $(p+1)/B^2$ precludes a
smaller pure denominator. In \eqref{eq:packet-shell} the scaled
inactive pure mass $E$ satisfies
$1-1/(p-1)-1/(6p)\le E\le1+1/(p-1)+1/(6p)<2$.
Every inactive label away from $B$ contributes at most $1/p^2$,
so exactly one occurs at $B$.

For $B=p$, Proposition~\ref{result.One-packet} gives the unique neutral
member $X$ and the asserted value on each active row by
\eqref{eq:packet-trace}. The vacuum Steinberg already uses all
$p-1$ active dimensions, so its complement is trivial: there $U_p$
is the identity, as are all its conjugates. Let $Q$ be the neutral
projector. The self-adjoint $Q(W_p-I)Q$ has rank one, since the
Steinberg neutral space is a line. For $w=W_pe_0$ its vacuum
entry is $w_0-1=-1/p$ and $w_X=\epsilon/p$ by the local column
formula. Hence $(W_p)_{YX}=w_Yw_X/(w_0-1)=-\epsilon w_Y$.
Local admissibility and Bruhat make the complementary Weyl act by $\epsilon$ on $X$, giving
$\chi_X(Y)=p s_{YX}=-p w_Y=-\epsilon a_Y$.
\end{proof}

\begin{lemma}\label{result:cross-weyl}
At a prime denominator $p$, let $a_Y$ be its proper coefficients.
For any nonunit row $X$ with $p\nmid b_X$,
$\sum a_Y\chi_X(Y)=\epsilon j$. If $m$ such rows of equal codegree
$f$ have all their proper hits in $Z$, then
$\sum_Z a_Y^2\ge mj^2/(f-m)$. With twelve denominator-four rows
and a denominator seven, or eleven such rows and exactly seven sevens,
there are six proper seven-free labels, with denominators $7k$,
$k\equiv\pm1\pmod7$, $k\ge6$.
\end{lemma}
\begin{proof}
At coprime denominators the nonunit entry of $s$ vanishes by lcm
integrality and strict unitary modulus at every embedding.
$W_pe_X=e_X$ by admissibility and Bruhat. Expanding
$\langle W_pe_0,se_X\rangle=1/b_X$ gives the identity.
Project $\sum_Za_Ye_Y$ and $\epsilon j e_0$ off the $m$ normalized
rows and apply Cauchy--Schwarz. Equality forces
$\sum_Za_Yd_Y=-\epsilon jm/(f-m)$ by taking the degree component.
Thus at twelve fours all four proper five-signs satisfy $4\mid b_Y$.
Three would be equality and the omitted signed degree would be
$\pm4$, impossible.

At twelve fours mass forces $j=1$ at seven. The norm six on at
least three degrees permits six signs or $2,1,1$ in magnitude.
In the latter case all three are singletons, hence their columns
are rational. Exactly one sign column $A$ has odd degree by the
signed degree sum. Call the other $B$. The four-row entries at $A$
are odd integers and the column norm forces signs. Their sum $w$
satisfies $|w|<16-m$ by \eqref{eq:saturation-gram}, so both signs
occur at least $m-7\ge4$ times. Column $B$ is even and at most
three entries can be nonzero. The double coefficient times its
column has fixed residue modulo four by dyadic principality, so
the cross identity modulo four forces a $B$-entry for every occurrence
of one sign at $A$, a contradiction. The congruence $ka_Y\equiv1$
gives the asserted indices.
\end{proof}

\begin{lemma}
If $M$ has prime support $2,t,5,7$, five and seven occur as denominators,
seven has one packet, and at least eight fours occur, their proper
supports have respectively four and six signs and meet in a singleton.
\end{lemma}
\begin{proof}
The local columns give $\sum_{H_5\cap H_7}a_{5,Y}a_{7,Y}=-1$, where
$H_p$ is the proper support. Two five-free witnesses omit two and $t$
and are outside $H_7$, as nonunit degrees are no prime powers. Hence
at most two labels are shared. With signs, the identity gives one.
If the seven-coefficients are $2,1,1$ in magnitude, their singleton
columns omit $2,t,5$, respectively in some order, and the shared
column $C$ has magnitude one. The double $Y$ has even degree, since
eight odd integer entries together with its seven-packet mass $4/7$
would exceed unit column norm. Let $A$ be the odd sign column.
There is a unique odd five-free witness $U$ by parity of its signed
degree sum and the two distinct even witnesses omitting $7,t$.

Multiply each of these four columns $V$ by $\epsilon\operatorname{sgn}a_{p,V}$
at its respective prime ($C$ at seven). By Lemma~\ref{result:packet-idempotent},
on columns $0,A,Y,U$ the seven- and five-packets give the rows
$v=(1,-1,-2,0)$ and $w=(1,0,0,-1)$.
On eight four-rows the $A,U$ entries are signs, and the even $Y$
column has at most one nonzero entry, of magnitude two, by the column
norm. Write $r_i=(1,a_i,y_i,u_i)$ for these rows, with $y_2=\cdots=y_8=0$.
Removing them and the packets gives
\[
 G=560I-80v^tv-112w^tw-35\sum_{i=1}^8r_i^tr_i\succeq0.
\]
There are $12\binom{10}3=1440$ choices of $r_1$ and the multiset of
the seven pairs $(a_i,u_i)$, $i>1$. Testing the fifteen principal
minors leaves only $r_1=(1,-1,2,-1)$ and multiplicities $(0,1,2,4)$
in lexicographic sign order. Then $G(6,-1,-7,1)^t=0$. Pairing with the vacuum row gives
$6-\tilde d_A-7\tilde d_Y+\tilde d_U=0$ for the signed degrees,
and the seven-column identity gives $\tilde d_A+2\tilde d_Y+\tilde d_C=1$.
Elimination contradicts $t\mid\tilde d_A,\tilde d_C,\tilde d_U$ and
$t\ne5$.
\end{proof}

\begin{lemma}\label{result:small-Gauss}
Write $T_b=\sum_{b_i=b}\theta_i$. For complete classes indexed by
$J$, of denominator lcm $L$, let $R$ be the complementary reciprocal
mass after deleting classes with $T_b=0$. Then
$L^2(R+1/M)<1$ implies $\sum_{b\in J}T_b/b^2=0$.
With denominator five present, eleven fours permit at most three
sixes. Twelve fours permit no six and at most two eights. In the
latter case the four proper five-free denominators are multiples of
twenty at least eighty. Two eights then give $T_8=0$ and preclude a
singleton nine.
\end{lemma}
\begin{proof}
Every conjugate of the integral sum $\sum(L/b)^2T_b$ is bounded by
$L^2(R+1/M)$ in the vacuum Gauss identity. Now $T_5=0$, $M\ge420$,
and eleven fours and four sixes leave $1/720$, forcing
$9T_4+4T_6=0$ contrary to reduction at two. With twelve fours the
remaining mass is $1/20$, giving $T_4=0$. A six would require
$1/M\ge2/36-1/20$. The indices $k_i$ of the five-signs are
multiples of four, congruent to the signs modulo five. Below sixteen
only $k_i=4$ with sign minus is possible, and even three terms of
size $1/16$ cannot offset $-1/4$ to $\pm1/n$. Three eights
would give $T_8=0$ with odd cardinality. With two eights the house
of $T_8$ is at most $64(3/160+1/420)<\sqrt2$, forcing zero by the
dyadic mean bound. A singleton nine then requires
$2/81\le3/160+1/M$.
\end{proof}
When $m_4$ is even, Lemmas~\ref{result:coupled-gap}
and~\ref{result:coupled-gap} give $T_4=T_6=0$ for
$m_6\le3,\ 36(R+1/M)<1$, or $m_6\le8,\ 72(R+1/M)<1$.
They give $T_4=0$ for $m_6\le5,\ m_8\in\{0,2\},\ 96(R+1/M)<1$.
Here $R$ omits exactly the classes under consideration together with
known zero sums.

\begin{lemma}\label{result:joint-projection}
Let $m$ rows have denominator four, $p$ be an odd prime dividing $M$,
and $U$ comprise the $p$-free nonunit hits of nonzero Galois mean
on these rows, or any larger Galois-stable set of $p$-free nonunits.
For positive orbit-constant integer weights $w_Y$ prime to $p$ put
$c_Y=w_Y[ad_Y/w_Y]_p$ on $U$, zero elsewhere, for $1\le a\le(p-1)/2$.
Brackets denote least absolute residues. The sums
$z_X=\sum c_Y\chi_X(Y)$ all have the same parity on the selected
rows. This parity is even if $2\mid w_Yd_Y$ for every $Y$. Let $l<0<u$ be the
neighboring integers of that parity congruent to $-a\pmod p$.
For $0\le j\le m$ integral put $S_j=(m-j)l+ju$,
$Q_j=(m-j)l^2+ju^2$. Then
\[
 B_a=\frac1{16}\min_j\left(Q_j+\frac{S_j^2}{16-m}\right)
 \le\|c\|^2,
 \qquad
 H_{16}(p;2,m):=
 \left\lceil\frac{\sum_a\lceil B_a\rceil}{\sum_a a^2}\right\rceil
 \le\sum_U w_Y^2 .
\]
With unknown parity take the minimum over both choices of $l,u$.
If $U$ contains every $p$-free nonunit and $w_Y=1$, a stronger
bound uses the piecewise-linear interpolant $L$ of the square on
$-a+p\ZZ$:
\begin{equation}\label{eq:full-census}
 \|c\|^2\ge\frac1{16}\min_j
       \{Q_j+(16-m)L(-S_j/(16-m))\}.
\end{equation}
\end{lemma}
\begin{proof}
Each $z_X$ is integral by Galois symmetry, $-a\pmod p$ by degree
orthogonality, and $\sum_Y c_Yd_Y\pmod2$ by dyadic principality.
Apply Bessel on these $m$ normalized rows to $c+te_0$ and optimize
over $t$, giving $16\|c\|^2\ge\sum z_X^2+(\sum z_X)^2/(16-m)$.
At fixed sum the least square cost uses adjacent lifts. If all
are on one side of zero they can be moved towards zero. Summing
over $a$ permutes the absolute residues at each $Y$, giving the
second bound. In the full case every row, including the vacuum,
has residue $-a$. With weights $1/f_X$ they have mean zero and
second moment $\|c\|^2$, so the remaining rows cost at least $L$ at
their mean. The neighboring lifts still suffice. Below $ml$ the
remaining mean is positive, where $L$ increases. Above $mu$ a
step of $2p$ costs at least $2p(2u+2p)$ in squares, exceeding
its saving since the slope of $L$ at negative mean is at most $p-2a$.
\end{proof}
In particular, taking weights one at even and two at odd $p$-free
nonunit hits yields $e_p+4o_p\ge H_{16}(p;2,m)$, where $e_p,o_p$
are the corresponding numbers of hits.

\subsection{The support of an even square dimension}
In this subsection $r\le44$, $M$ is even, and $m$ counts denominator-four
labels. A reservation $(h,b)$ charges $h$ distinct positions of denominator
at least $b$. After disjoint reservations $(h_i,b_i)$, with floor $c$
elsewhere, the upper mass is
\[
 U((h_i,b_i);c)=\sum h_i/b_i^2+(44-\sum h_i)/c^2.
\]
A negative remaining count or an upper mass below one excludes the
profile. Padding by unreserved positions covers every smaller rank.

\begin{lemma}[Three-prime support]\label{result.Three-prime-rank}
At least four primes divide $M$.
\end{lemma}
\begin{proof}
Suppose the primes of $M$ are $2<p<t$. Each occurs to
exponent at least two. A $u$-free witness in $P_v$ ($u\ne v$)
has both other primes in its degree and $v$ in its denominator.
Thus $M\ge900$ and no prime denominator occurs. Write $m\le13$ for
the multiplicity of four. Other nonunit denominators have floor six
(eight if $p>3$). We use the following position counts to bound the
sum of reciprocal squares, padding to $44$ places.

Put $A=4$ if $v_2(M)=2$, and $A=8$ otherwise. The first case has
$m\le3$. Odd degrees occupy $3\bmod4$ nonunit places. These include
the $m$ fours when $A=4$, and one further denominator divisible by
$At$ from $P_t$. The other odd places have floor $4p$ if $A=4$,
and $8$ otherwise. For $s=p,t$ let $n_s$ count the $s$-free
nonunits. At least $k_s=H_{16}(s;2,m)$ have denominators divisible
by $4s^{v_s(M)}$ if $m\ge4$. Put $k_s=0$ otherwise. Use
$n_s\ge\max(2,k_s)$, $n_3\equiv2\pmod3$, $n_5\ge3$, and the
stronger bound \eqref{eq:full-census} at five and seven when $m\ge4$.
The three omission sets are disjoint. When $m<4$ one $p$-free
denominator is divisible by $p^{v_p(M)}t$ and one $t$-free denominator
by $t^{v_t(M)}p$, by principal witnesses.

If $t$ mixes with two or $p$, reserve respectively $4(t-1)$ or
$p(t-1)$ active multiples of $2t$ or $pt$. These avoid the four-row
hits and, when $m>0$, the odd degrees, whose twists are dyadic. For
the odd interaction with $m\le3$ simply reserve the entire $pt$
places of the spherical summand, allowing the vacuum among them.
If $t$ is isolated, first test $t(t-1)$ pure places of floor $t^2$,
allowing them to supply $t$-free places other than four-row hits or
the principal-$p$ witness. Below this count sparse height gives
$M_t=t^{2c-1}$, $c\ge2$, and every nontrivial $t$-part has size
at least $t^2$. Nonminimal activity then reserves $2(t-1)$ pure
$t$-divisible degrees, by primitive balance or a higher orbit.
With fewer than $t(t-1)$ pure labels, minimal activity leaves fewer
than $(t-1)^2$ inactive pure labels. Thus Lemma~\ref{result:packet-idempotent} gives $t$
such degrees at the common denominator $t^c$.
Adding the reciprocal squares of these floors leaves only
$(p,t)=(3,5)$ with minimal activity, $A=8$, and
$m\in\{4,5,7,8,9,10,11,12,13\}$.

Write $M=2^a3^b5^{2c-1}$. The vacuum Weyl congruence gives
$2^a3^b\equiv\pm1\pmod{125}$. Truncate $a,b,c$ at the first
exponents with respective powers exceeding $K=65\,536$, the
last value representing its whole tail. Reserve the vacuum, the
$m$ fours, the five places at $5^c$ and the minimum omission
counts above. For each remaining position enumerate the permitted
divisors $d\le K$ of $M$ with their exact reciprocal squares.
A nonunit degree omits at most one prime. A larger denominator of
floor $F$ contributes in $[0,\max(K+1,F)^{-2}]$, also allowing
padding. The vacuum has positive mass at most $(2^a3^b5^{2c-1})^{-2}$ at the chosen lower endpoints.
If a valuation was truncated, omit its congruence test when necessary
and use the lower endpoint zero for a truncated packet mass.
Enumerate in order within each group, discarding intervals that
cannot contain one. The nine profiles leave no intervals after
$18\,095$ nodes (\texttt{three\_prime\_cutoff.py}).
\end{proof}

\begin{lemma}[Restrictions on mixed twists]
No twist is active at two odd primes. An odd prime mixed with two
has at least $4(p-1)$ active multiples of $2p$, and $p\le5$.
\end{lemma}
\begin{proof}
Suppose two odd primes mix. The spherical summand and the activity
reservations exclude a product $K\ge33$. At $K=21$ the third odd
prime is five. A fourth odd prime gives mass at most
$11/16+1/15+1/7+1/11<1$. Its summand is disjoint as their product
would exceed the rank. A second odd edge, or dyadic interaction at
the third prime, is excluded by the counts and mass (for edges $15,21$
use $12/225+18/441+14/16<1$. For $21$ and dyadic five use
$16/100+18/441+10/16<1$). At $K=21$ reserve its $22$ mixed labels.
Absence of five leaves mass below one even with thirteen fours.
Its presence reserves five labels and four proper witnesses. At
most eleven fours have insufficient mass, and thirteen overshoot.
At twelve all positions are reserved. The four proper five-free
denominators are at least $80$ by Lemma~\ref{result:small-Gauss}.
If any of the $22$ multiples of $21$ exceeds $21$, then
\[
 1\le12/16+1/5+21/21^2+1/42^2+4/80^2+1/420^2<1.
\]
Thus all equal $21$. At equality in Corollary~\ref{result:fiber-squares}
there are exactly two nontrivial ternary fibers meeting seven-activity,
with eight eligible labels apiece. A higher ternary phase would require
three such fibers, a higher seven-phase $21$ labels in one fiber.
Each therefore has six singleton primitive seven-phases and two eligible
seven-neutral labels, all of the same weight. The fiber moments give
$(1-a_1)^2+(1-a_2)^2=8$ with $a_1,a_2\in7\ZZ$, impossible.

At $K=15$ the third odd prime $u$ is isolated. A mixed reservation
and the eighteen multiples of $15$ leave insufficient mass. Thus
$u<44$ by the $u$ pure places. Reserve the vacuum, $m$ fours, these
eighteen places, and either $u$ at $u^2$ or $uj$ at $u$. In the latter
case partition $j(u-j)$ into at least three positive squares $a_i^2$.
The proper witness floors are $uk_i$, with $k_i>1$ least satisfying
$a_ik_i\equiv\pm1\pmod u$. Use Lemma~\ref{result:cross-weyl} at seven.
For $m\ge7$, Lemma~\ref{result:saturation} gives $8\mid M$: reserve three
odd degrees, imposing $8\mid k_i$ on their overlap. The proper
witnesses include an odd degree by gcd one and an even degree since
every odd $u$-free degree contains three and five.
At $j=1$ the prime twist sum vanishes, so odd $m$ requires
$16(1-m/16-1/u+1/(30u))\ge1$. Every remaining upper mass is below one
(\texttt{arithmetic/isolation44.py}).

A dyadic interaction at $p\ge7$ claims at least twenty-four positions
of floor fourteen. A five-packet, its four proper witnesses and two
ternary active labels then give insufficient mass. Otherwise six
positions are active at the other odd primes. For $m\le10$ the bound
is $10/16+24/196+9/36+1/210^2<1$. For $m\ge11$ three odd nonunits
require at least $1+24+6+11+3=45$ places. The count at a mixed prime
is Corollary~\ref{result:mixed}.
\end{proof}

\begin{lemma}
If every active twist has prime-power order and $\xi=\pm1$, the
deficit $1-\xi/M=\sum(1-\langle\theta_i\rangle)/b_i^2$ is a sum of
digits. At odd $p$ a digit $p^{1-2k}$ has $p^k\mid M$ and costs at
least $p-1$ active labels. At a prime denominator the full contribution
is $j/p$. At two each active label contributes $2^{-2k}$ or $2^{1-2k}$.
\end{lemma}
\begin{proof}
Local splitting makes an active denominator a power of its twist prime.
The primitive $p$-phases have equal weight $W$ and deficit $pW$. Use
the least numerous fiber to charge $p-1$ labels per term. Higher
phases have mean zero and occur in equal-degree packets divisible
by $p$. The dyadic formula follows from the root trace.
\end{proof}
In the resulting exact recursion, if the current nonzero remainder
is $A$, at most $s$ digits remain, and their largest size is $1/u$,
then $|A|\le s/u+u^{-1/2}$, since the final error is $\pm1/M$.
Each choice imposes its displayed divisor on $M$. A zero intermediate remainder would leave a positive sum of digits
equal to $1/M$. At an odd $p\mid M$ the sum of the remaining $p$-digits
would then have valuation $-v_p(M)$, since every other digit is
$p$-integral. Its reduced denominator would be $M_p$, giving a sum
at least $1/M_p>1/M$.

\begin{lemma}
Exactly four primes divide $M$.
\end{lemma}
\begin{proof}
Each odd $p$ needs $p$ pure labels if isolated and $4(p-1)$ if mixed.
Five odd primes require at least $39$ places and only the prime three could
mix with two. The upper mass with all isolated is
$1/27+1/5+1/7+1/11+1/13+4/16+1/2310^2<1$.
Mixing three would require $8+5+7+11+13+1=45$ labels. Larger choices
only reduce the bound.

For four odd primes reserve at each $p$ one of $(p,p^2)$,
$(4(p-1),2p)$, or $(p,p)$ according as $p$ is absent and isolated,
mixed, or present. A present denominator contributes at least three
proper witnesses at $2p$, or four at $20$ if $p=5$. For $m\ge9$
reserve three odd labels of floor eight, allowing overlap with the
witnesses by taking coordinatewise maxima of their decreasing floors.
Use $M\ge2310$ and default floor six. Enumeration gives $49$
profiles, of which the $38$ all-primary ones are excluded by the
preceding deficit recursion. Here at least one prime denominator
is present. Subtract its $j/p$ and $(m-\langle T_4\rangle)/16$,
and charge the vacuum, fours, prime classes and largest proper
witness quota before enumerating the remaining digits.

In the other eleven profiles eight reserved denominators are multiples
of six supported on two and three. Reserve also $p$ pure powers
for each absent isolated prime and $p$ copies of each present prime.
At most eight positions remain. A present prime has one packet.
This holds at five, and another packet at $p\ge7$ would leave fewer
than the three required proper degrees. On the remaining positions
its Weyl coefficients satisfy
\[
 \sum_i a_{pi}^2=p-1,\qquad \sum_i a_{pi}/b_i=\epsilon/M,\qquad
 \sum_i a_{pi}a_{ti}=-1\quad(p\ne t).
\]
Indeed $(W_p-I)e_0$ and $(W_t-I)e_0$ are orthogonal since no summand
is nontrivial at both primes. Deleting the vacuum gives the last
identity. If $p\nmid b_i$ then $a_{pi}=0$. Otherwise it is the least
absolute inverse of $b_i/p$ modulo $p$.

Enumerate the coefficient magnitudes with the indicated norms, at
least three nonzero entries per row and odd pairwise dot products.
Within each resulting column type try the compatible denominators
through $256$ in increasing order, allowing repetitions. A larger
denominator contributes mass in $[0,257^{-2}]$ and at most
$|a_{pi}|/257$ in modulus to the signed sum. The zero endpoints
allow padding to rank $44$. For the vacuum use
$M\ge4(\prod_{p\mid M,\,p>2}p)\prod_{t}t$, where $t$ runs over
the absent isolated primes. The mass test leaves eighteen intervals. The signed sums exclude all eighteen (\texttt{support\_five.py}).
\end{proof}

\begin{proposition}[Multiplicity at denominator four]\label{result:thirteen-34}
With $m\ge9$ a twist is active at two and three, and every other
odd prime is isolated. At $m=9,10,11,12,13$ the numbers of admissible
mixed profiles are $12,11,9,6,1$. Through rank $41$, $m\le12$.
Through $40$ the odd primes are $3,5,7$, with five and seven
isolated. For $m=9,10,11$ in the latter range, absence of denominator
five forces seven. At eleven both cannot occur, and at twelve seven
is absent.
\end{proposition}
\begin{proof}
Here $M/4$ is even and $m\le13$. The vacuum, $m$ fours and at least
three odd hits avoid all odd activity. At an isolated absent $p$ reserve
$(p,p^2)$, at a mixed $p$ reserve $(4(p-1),2p)$, and at a present
$p$ reserve its prime packet. For every odd $p$ record a divisor $P_p$
of the denominator of a $p$-free degree: use $p$ for a mixed or present
$p$, $9$ for absent isolated three, $p^3$ for other absent isolated
primes. In the last case the sparse-height bound applies. Indeed the
vacuum, fours, odd labels and other prime activity leave fewer than
$p(p-1)$ positions at $p\ge7$. At absent isolated five, twenty active
labels give mass at most $20/625+10/16+3/64+7/49+3/36+1/210^2<1$.
At eighteen, three further five-free positions have floor $125$
since sparse height applies below twenty. Eleven fours still give
insufficient mass. Twelve require precisely three remaining positions,
three pure ternary and seven pure seven labels on $2,3,5,7$. Those
three remaining would be both five-free and odd. Without seven the
mass is too small, while its packet needs an odd seven-free witness,
necessarily then of prime-power degree. Thirteen fours cannot fit.
Hence at most sixteen are five-active. An odd prime $\ge23$ exceeds
the rank or has insufficient mass already at $23$.

Group the remaining labels by the odd primes $J$ omitted by their
degree. A denominator floor is $c\prod_{p\in J}P_p$, with $c=8,4,1$
for an odd degree, even four-row hit, and other even degree. The
possible $|J|$ are at most one for odd degrees and two otherwise.
Impose the weighted and unweighted bounds of Lemma~\ref{result:joint-projection}
on the hits (in the latter allow either common parity of their row
sums). At a present prime the proper free degrees omit each other
prime. Of the four five-witnesses exactly one is odd, by parity and
the two distinct even witnesses omitting the other odd primes.
At least $\lceil m/(16-m)\rceil$ are hits by Lemma~\ref{result:cross-weyl}.
At twelve use also Lemmas~\ref{result:cross-weyl} and~\ref{result:small-Gauss}.
Thirteen fours exclude five by mass and seven by the impossibility
of $T_4=0$ with odd cardinality (one packet leaves $5/112$, two
overshoot). Set $M\ge8\prod P_p$ in the vacuum slot.
These are the finite omission tests. Retain every nonnegative
profile for which the mass can reach one.

The omission sums give $41$ mixed signatures. Only two mix five, both
on $2,3,5,7$ with three isolated and seven present. Their omission
lists at ten and eleven fours number five and one. Each contains three
odd witnesses omitting $3,5,7$, respectively, and two even hits
omitting $\{3,7\},\{5,7\}$. At eleven the sole other witness omits
$\{3,5\}$. All positions are accounted for, with seven sevens but only
three proper seven-free labels, contrary to Lemma~\ref{result:cross-weyl}.
At ten, after the vacuum, fours, sixteen mixed labels, seven sevens and
the five witnesses, at most five positions can be ternary-active.
Lemma~\ref{result:activity-height-gap} gives $27\mid M$ and excludes
six. Reserving three pure ternary labels gives the contradiction
\[
 1\le\frac{10}{16}+\frac{16}{100}+\frac3{81}+\frac17+\frac2{64}
 +\sum_{b\in\{216,40,56,756,140,7560\}} b^{-2}<1.
\]

For the all-primary case reserve at least three neutral labels by the
four singleton twist components (Corollary~\ref{result:mixed-neutral-rank}).
Take the maximum with the proper-witness quota. Five and another prime
packet require at least six proper labels, since two at the latter
must omit two and the remaining odd prime separately and cannot be
five-free. Subtract prime contributions $j/p$ and $(m-t)/16$,
$t=\langle T_4\rangle$. The digit recursion with these reservations
leaves, from $46$ profiles, just $M=136080,15482880$, both with nine
fours and one packet at five and seven (a second seven-packet
overshoots). Their proper signed unions number $13,39$. Both prime
twist sums vanish and the residual mass gives $|t|\le1$. Enumerate
the ternary digits $3/b^2$ at $b=9,27,\ldots\mid M$, each costing
two places, and all dyadic splittings. This leaves $0,218$ profiles.
For ternary deficit $G$, $h$ open positions and $v$ nines,
$2G/3\le v/81+(h-v)/729$. Beyond the nines, all remaining degrees
are multiples of $35$ dividing $M$. Divide the residual square sum
by $35^2$ and test all divisor squares with zero padding. None of the
$5092$ targets on at most five slots occurs. The omission sums,
deficit recursion and terminal square sums are the finite comparisons
implemented by \texttt{preflight.py} in the census source.
Using the actual rank in the same sums gives the asserted smaller-rank
restrictions (ten mixed profiles through rank forty).
\end{proof}

\section{The denominator census}\label{sec:verify}

Assume henceforth that $r\le44$. The dimension is even by Lemma~\ref{general:odd-small} and square-full by Lemma~\ref{lem:basic-table}, so $2\mid M$.
Let $B$ be the least nonunit denominator, of multiplicity $m$, and put
$n=M/B$, $f=qB^2$. The ordered denominators determine $M$, their lcm,
since the nonunit degrees have gcd one. We have $f<r$. Equality
$f=r-1$ would give $\sum_{i\ne0}(n^2-d_i^2)=1$, whereas each
positive summand is at least $2n-1$. The small-codegree bounds thus leave
$16,25,36$ for square $D$ and $12,18,20,24,27,28,32,40$ otherwise.
Apply Lemmas~\ref{result:small-effective} and~\ref{block:packets}, and
$m\le f-1-(p-1)/2$ for odd $p\Vert f$
\cite[Proposition~19.9]{BDPcore}, by principal height and the primary mass identities.

\paragraph{The search interval.}
Use \eqref{eq:interval}, with $h$ nonunit positions still empty.
For the two-degree terminal equation it is enough to try
$k^2\le(1+2L^2)/E$ and the divisors of the prime-to-$k$ part of $L$.
With three degrees $d,e,g$ left and $k$ even,
$Ek^2=1+d^2+e^2+g^2$ modulo four and eight makes all three odd,
$v_2(k)=1$, and $E$ odd. No new prime can enter $M$ at the last
two positions, by Lemma~\ref{block:witness}. 

\paragraph{Reserving positions.}
If $v$ of the $h$ missing nonunits have respective denominator floors
$a_i$ and all denominators are at least $b$, the residual mass is at most
$\sum a_i^{-2}+(h-v)b^{-2}+M^{-2}$. Combine overlapping reservations
by taking coordinatewise maxima after sorting. For
$W=\sum\lceil L^2/a_i^2\rceil$ and $E>W+1$ this gives
$b^2\le(h-v)L^2/(E-W-1)$, since $L\mid M$. The activity bound on
every set of odd primes $p\mid M$, $p\nmid q$, lets Hall's theorem reserve
disjoint groups of $B(q)(p-1)$ labels for these primes. A closed class of
size $s$ supplies at most $h_p\lfloor s/h_p\rfloor$ at $p$,
$h_p=(p-1)/2$. Credit an open class in full. The support contains the
primes of $2qB$ and satisfies $B(q)(1+\sum_{p\mid M,\,p\nmid2q}(p-1))\le r$.

For $q=3$ the neutral fiber at the unramified odd primes costs four
labels in place of $B(3)=3$. Its principal ternary block contains the
vacuum. A phase of order at least nine supplies three more labels.
Otherwise the order-three mass is twice the inactive mass. With at most
three labels, all nonunit degrees exceeding two, this requires
$b^2=2(1+a^2)$ for nonunit degrees $a,b$. Modulo three they are
three-free and modulo four $a$ is odd, so $(a,b)=1$. Their $S$-entry
vanishes, contrary to principality at three.

For $q=3,5$, $p\ge5$ distinct from $q$, put $F=(q+3)/2$.
Every nonempty fiber away from $q$ costs $F$ labels, even with dyadic
phase fixed, as the complementary projectors commute with the $q$-factor.
With fewer than $2F(p-1)$ active at $p$, their phases have order $p$:
order at least $p^2$ costs $Fph_p$ labels. The fiber with trivial
complementary phases contains an active $p$-phase, since otherwise the
local unipotents would preserve the vacuum line. Equal primitive masses
give all $p-1$ roots there. If an active $X$ has $t\mid f_X$ outside
$\{p,q\}$, its fiber away from $t$ must contain two $t$-phases by the
same invariant-line argument. Thus some nontrivial $t$-phase also contains
all $p-1$ roots, costing $2F(p-1)$ labels. Hence active denominators are
supported on $p,q$. At $q$-free denominators their defect-one blocks must
have equal degrees, since a mixed block includes $z,2z$. Their activity
multiplicity is consequently divisible by $qh_p$.

Tile $p$-free denominators for $p\mid q$. At two the counts are even.
At three and five, after obligations from $z/2$, a mixed block uses
respectively one or two copies of $z$ and two copies of $2z$.
The least mixed count permitting equal blocks at $z$ suffices: $p$
mixed blocks may be exchanged for equal blocks at both levels.

At $q=3$ on support $\{2,3,p\}$, any position at denominator $2$ or
$p$ needs two at twice its denominator. An equal ternary block there
would require a fourth prime by Proposition~\ref{result:effective}.
At degree $d$, constituents of $XX^*-1$ prime to each $p\mid d$
have total degree congruent to $-1$ modulo the least $p$-part of a
$p$-divisible degree. On three-prime support a nonunit can omit only
one prime. Thus for the three least nonunit degrees $d\le e\le a$,
with repetition, $e+a\le d^2-1$, since at least two distinct
constituents other than $X$ are required.

For $p\ne t$ in the support, a nonunit in $P_p$ omits $t$, giving a
denominator multiple of $pt^{v_t(M)}$. Orbit deletion forces two
$t$-free nonunits, and reduction of $D$ modulo four or five forces
three for $t=2,5$. With three support primes their omissions are
disjoint, every $t^2\mid M$, and the two- and three-free counts are
$3\bmod4$ and $2\bmod3$, respectively. At an undecided height take
the weaker bound for a fixed or increased valuation. With four
primes a nonunit omits at most two. Distinct omitted pairs require
distinct positions. The proper support at a prime denominator must
omit each other prime somewhere by its gcd-one condition.

\paragraph{Square dimension.}
There are four primes, and the only possible edges are $2$--$3$ and
$2$--$5$ (Lemma~\ref{result.Three-prime-rank} and the subsequent support
bounds). Use the multiplicity restrictions of Proposition~\ref{result:thirteen-34}. At an isolated odd $p$ reserve its $p$
pure nonunit positions. At $p\ge7$ there are at most $37$ active labels
after reserving the vacuum and the two other odd activities. At seven
this number is a multiple of three, hence at most $36$. At isolated
five with no denominator five there are at most sixteen. For at least
nine fours this is the high-four bound, and with at most eight fours
eighteen active labels would give total mass at most
\[
 18/625+6/49+2/36+8/16+9/36+1/210^2<1 .
\]
The count is even. Thus at every isolated $p\ge5$ the sparse-height
inequality applies. In the absence of denominator $p$, no mixed
denominator has $p$-part $p$.
At the greatest active level $p^a$, $v_p(M)=2a-1$ and at least
$p-1$ labels occur there (a higher-order packet suffices, and for
primitive packets reduce equality of the phase masses modulo $p^2$
after scaling by $p^{2a}$). These give denominator floors both for
the active labels and the omitted-prime witnesses. At isolated three reserve only two equal pure denominators, by the same
primitive-mass reduction or a higher-order orbit.
If seven occurs, five does not, and $m_4+m_6+27>r$, five cannot be
mixed: it would need sixteen active labels in addition to the vacuum,
the fours and sixes, and either fourteen sevens or seven sevens and
three proper, twist-trivial witnesses.

For an isolated $p$, let $E_p$ be the mass of all pure $p$-power
nonunits and $C$ the sum, over observed mixed $b=uk$ with $u$ its
$p$-part, of $\min_{ak\equiv1\, (u)}(a/u)^2$ ($u>1$). Since the pure
entries use $E_p$ of the squared norm $2g_p-g_p^2$ and
$0<g_p\le pE_p/(p-1)$,
\[
 C<(p+1)E_p/(p-1).
\]
With fewer than $p^2$ active labels, equality of the primitive masses
at successive levels gives equal base-$p^2$ digits at each pure
denominator. Higher-order orbits contribute in multiples of $p$
and have size at least $ph_p$. Thus each class of $n_b$ labels contributes
$pc_b/b^2$ to $g_p$ with $(p-1)c_b\le n_b$. For $H$ remaining
pure positions of floor $b_*$ it follows that
\[
 C\le\sum_{b\ {\rm closed\ pure}}
 \frac{2p\lfloor n_b/(p-1)\rfloor-n_b}{b^2}
 +\frac{(p+1)\lfloor H/(p-1)\rfloor}{b_*^2}.
\]
Indeed $2pc_b$ norm units use at least $(p-1)c_b$ units on the pure
positions themselves. If a mixed dyadic--$p$ phase occurs,
each of its required fibers uses four positions of one positive dyadic
height. In the ternary case the two primitive masses agree. A partial
difference at denominators of lcm $W$ is a multiple of $1/W^2$. If the
remaining eligible mass is smaller, the difference is zero.
Between two such cuts a singleton at maximal dyadic or ternary
denominator valuation must have coefficient zero, by reduction at that
prime. A singleton also cannot carry an order-nine phase. Removing these
positions before reserving the quartets preserves every solution.

At a prime denominator $p$ of multiplicity $pj$, use the signed
coefficient equations of Proposition~\ref{result.One-packet} on its proper
witnesses. For $p=7$ the budget
$j(7-j)\le12$ forces each coefficient to be the least absolute
inverse of $b/7$ modulo seven. For any partial list, with remaining
budget $K$, its missing sum $\sum a/(b/p)$ is bounded in either
sign by $K/\min |a|(b/p)$, minimizing over admissible remaining
terms of that sign. With one packet at $p=5,7$ write $W$ for the
lcm of a nonempty proper initial witness list, $s=\sum aW/b$.
Then $s\equiv (W/p)\sum a^2\not\equiv0\pmod p$ since the total
budget is $p-1$. For the next denominator $z$, if the missing
absolute coefficients sum to $K$, the signed equation gives
$z\le(K+1/6)W/|s|$. At the last witness put $d=W/\gcd(W,z)$,
$t=z/\gcd(W,z)$. The gcd-one condition gives $M=Wt$ and
$st+a_zd=\pm1$. Testing $d\mid W$ is exhaustive.
For simultaneous single packets at five and seven at most two proper
labels are shared. In each family separate labels must omit each of
the two other primes. In the three-label seven family at most one
can be shared. The cross-column sum of products on the shared
labels is $-1$, so exactly one is shared, with opposite signs,
also when the seven coefficients are all signs. Its denominator
$35k$ has $k\equiv\pm17\pmod{35}$. The remaining seven witnesses
are obtained from their same signed equation, with the common sign
specified by the vacuum congruence. With a present five, a mixed
ternary phase and $m_4\ge9$, the large isolated prime has at most
$r-m_4-20$ pure positions, after the vacuum, fours, five-packet,
four proper witnesses, two further odd degrees and eight mixed positions.
When this is at most $2(p-1)$ and $p$ itself is absent, the minimal
packet bound gives $p$ positions at the first pure level $B=p^c$ and
$M_p=B^2/p$, with $M/M_p\equiv-\epsilon\pmod{M_p}$.

\paragraph{Gauss sums.}
For a complete class write $T_b=\sum_{b_i=b}\theta_i$.
If selected classes have denominator lcm $W$ and the complementary
mass is at most $R$ after discarding known zero sums, every conjugate
of $W^2\sum T_b/b^2$ has modulus at most $W^2(R+1/M)$.
If this bound is below one, integrality forces the sum to vanish. Reduction modulo
a prime detects nonvanishing when a unique maximal-height class has
one label, or is a pure $p$-class with multiplicity prime to $p$.
One may repeat after each zero sum, or apply the dyadic cancellation
bound to two classes. A selection of individual terms suffices as well. After
a zero seven-sum, add one twelfth-denominator twist to $9T_4+4T_6$
when $m_4$ is even. Reduction at two then leaves a single unit.
When $m_4$ is odd the four-class alone gives that unit. This proves the corresponding suffix tests by
using an upper bound for $R$. For example, with ten fours, five fives
and six sixes the five-sum vanishes. The short Gauss relation with
$72(1/120+1/840)<1$ gives $T_4=T_6=0$. The remaining mass is
$1/120$. A fourth prime $p\ge7$ is isolated and absent as a denominator
since $p-1\le41$ and a prime packet alone would exceed that mass.
Thus $M\ge8\cdot3\cdot5\cdot7^3$.

For an isolated odd prime $p$, put
$M_p=p^{v_p(M)}$. After subtracting the deficits
$(1-\langle\theta_i\rangle)/b_i^2$ of any complete nonunit classes
from one, the remainder is nonzero. Otherwise the unassigned
nonnegative deficits sum to $1/M$ by the Gauss identity.
Their $p$-active part $P$ has $2P$ of $p$-power denominator. Only
pure $p$-denominators contribute, and every class sum is invariant
under square multipliers, whose index is two at $p$. All other
nonzero terms have denominator prime to $p$, since all classes are
square-multiplier invariant and their Galois means have
power-of-two denominator. Hence the denominator of $2P$ divides
$M_p$, giving either $P=0$ or $P\ge1/(2M_p)>1/M$.
The former contradicts the denominator of $1/M$. The unassigned
deficits are bounded by twice their mass, with terminal error $1/M$.
At isolated $p$ with fewer than $p^2$ active labels one uses
the digits $pc_b$ above. At all other classes the root means are
$\mu(k)/\varphi(k)$.

\paragraph{Excluding a complete list.}
Apply five tests in order. First, if $d$ is the least nonunit degree and
$E_p=\sum_{i\ne0,\,p\mid f_i}d_i^2$, Lemma~\ref{block:packets} requires
$2pdE_p\ge(p-1)D$ for every $p\mid D$.

Second, at a prime denominator $p\nmid q$, put $n=M/p$. For $q=1$ this
gives a row of codegree $p^2$. For $q=2$, or for $q=3,5$ when neither
$2n$ nor $n/2$ occurs, the defect-one block of the ramified prime has
equal degrees and gives an effective row of codegree $p^2$.
Proposition~\ref{result:effective} then requires gcd one
among the nonunit degrees prime to $p$.
Third, tile the defect-one classes as above.

Fourth, let $A_p$ consist of the nonunit degrees $y$ with $p\mid f_y$ such
that, for every $p$-free nonunit degree $x$ of multiplicity $m$,
\[
 g=(x,y)>1,\qquad 2mg>p,\qquad m\,\lcm(x,y)^2<D .
\]
It contains every nonunit degree of $P_p$, so $\gcd A_p=1$.
Indeed a principal row meets the entire $x$-class, giving $g>1$.
Its entries divided by $x/g$ are nonzero algebraic integers.
Their sum is an integer congruent to $mg$ modulo $p$ and strictly
smaller than $mg$ in modulus. Their mean squared moduli are at
least one. The vacuum uses one further unit of row norm.

The remaining lists have $q=1$ and a denominator five. By
Lemma~\ref{result:five-reduction} there are four proper five-free
nonunit labels, i.e. those of degree $d\ne M/5$ prime to five.
On these an active five-row has signs $s_d\equiv d\pmod5$
by \eqref{eq:packet-trace} and principality. It vanishes on
five-divisible degrees. Its own degree class lies outside $P_3$
by principal height. Thus some subset of the four labels with
$d\in A_3$ must satisfy $1+\sum s_d d\equiv0\pmod{M_3}$
by \eqref{eq:principal-trace}.

\begin{proof}[Conclusion of the proof of Theorem~\ref{alg:rank}]
The denominator recurrence gives the following numbers of complete lists.
\[
\begin{array}{c|rrrrrrr|r}
 r&1\text{--}38&39&40&41&42&43&44&\text{total}\\ \hline
 q=1&0&1&2&17&111&914&6269&7314\\
 q>1&0&0&0&4&65&723&10093&10885
\end{array}
\]
Of these $18\,199$ lists, the five tests leave successively
$8\,207$, $6\,194$, $6\,166$, $244$ and $0$. The enumeration covers
$12\,273$ tasks. At a fixed-width arithmetic limit the unfinished branch
continues in arbitrary precision. There are $41\,915$ such continuations.
The supplement records each list and its first exclusion. Running
\texttt{python3 verify.py --full} regenerates the census. Every admissible
list follows the ordered recurrence and satisfies its pruning inequalities.
Hence $r\ge45$.
\end{proof}

\subsection*{AI use declaration}
GPT-6 Astra and GPT-6.2 Astra assisted with the development and drafting of proofs, the literature search, and the implementation and checking of the finite calculations.

\begin{thebibliography}{99}
\bibitem{ABPP}
M.~A.~Alekseyev, W.~Bruns, S.~Palcoux and F.~V.~Petrov,
\href{https://doi.org/10.1063/5.0334158}{\emph{Classification of integral
modular data up to rank 13}}, J. Math. Phys. \textbf{67} (2026), no.~8, 081701. Numbering follows the \href{https://arxiv.org/abs/2302.01613v10}{open author version}.

\bibitem{BDPcore}
S.~Burciu, J.~Dong and S.~Palcoux,
\href{https://hexagonmath.org/2611.00233v3}{\emph{Prime factors of fusion-category dimensions: weak group-theoreticity, rank and decomposition}},
Hexagon preprint 2611.00233v3 (2026).

\bibitem{BPBurnside}
S.~Burciu and S.~Palcoux,
\href{https://doi.org/10.1142/S0219199726500070}{\emph{Burnside-type
results for fusion categories}}, Commun. Contemp. Math. (2025), article 2650007.

\bibitem{CMSCyclotomic}
F.~Calegari, S.~Morrison and N.~Snyder,
\href{https://doi.org/10.1007/s00220-010-1136-2}{\emph{Cyclotomic integers, fusion categories, and subfactors}},
Comm. Math. Phys. \textbf{303} (2011), 845--896.

\bibitem{DLNCongruence}
C.~Dong, X.~Lin, and S.-H.~Ng,
\href{https://doi.org/10.2140/ant.2015.9.2121}{\emph{Congruence property in conformal field theory}},
Algebra \& Number Theory \textbf{9} (2015), no.~9, 2121--2166.

\bibitem{ENOannals}
P.~Etingof, D.~Nikshych, and V.~Ostrik,
\href{https://doi.org/10.4007/annals.2005.162.581}{\emph{On fusion categories}},
Ann. of Math. \textbf{162} (2005), no.~2, 581--642.

\bibitem{HR}
S.-M.~Hong and E.~C.~Rowell,
\href{https://doi.org/10.1016/j.jalgebra.2009.11.044}{\emph{On the classification of the Grothendieck rings of non-self-dual modular categories}},
J. Algebra \textbf{324} (2010), no.~5, 1000--1015.

\bibitem{LPR}
Z.~Liu, S.~Palcoux, and Y.~Ren,
\href{https://doi.org/10.1142/S0129167X23500301}{\emph{Interpolated family of non-group-like simple integral fusion rings of Lie type}},
Internat. J. Math. \textbf{34} (2023), no.~6, 2350030.

\bibitem{Nevins}
M.~Nevins, \href{https://doi.org/10.4153/CJM-2005-026-1}{\emph{Branching
rules for principal series representations of $\mathrm{SL}(2)$ over a
$p$-adic field}}, Canad. J. Math. \textbf{57} (2005), 648--672.

\bibitem{NSind}
S.-H. Ng and P. Schauenburg, \href{https://doi.org/10.1007/s00220-010-1096-6}{\emph{Congruence subgroups and generalized Frobenius--Schur indicators}}, Comm. Math. Phys. \textbf{300} (2010), 1--46.

\bibitem{PrasadSL}
A.~Prasad, \href{https://arxiv.org/abs/0712.4051}{\emph{Representations
of $\mathrm{GL}_2(\mathbb F_q)$ and $\mathrm{SL}_2(\mathbb F_q)$, and
some remarks about $\mathrm{GL}_n(\mathbb F_q)$}}, 2007.

\bibitem{Stacks}
The Stacks Project Authors, \emph{The Stacks Project},
\href{https://stacks.math.columbia.edu/tag/04GH}{Tag 04GH}.
\end{thebibliography}
\end{document}
